\label{chap:83-01}

% RESULTADO_RECUPERADO. Owner integro: sections/hmt/03_app_ortograma.tex
% Destino causal: capitulo 1; la jerarquia se subordina sin alterar el owner.
% Los dos bloques científicos del owner se activan en una única residencia.
\begingroup
\def\HMTAPPFormal{1}
\def\HMTAPPDevelopments{1}
\begin{HMTScientificOwnerSubordinated}
\input{sections/hmt/03_app_ortograma}
\end{HMTScientificOwnerSubordinated}
\endgroup

% Unidad U003. Redacción autónoma; tablas y censos reconstruidos desde 1063/live.

\section{Aritmetic Matrix of APP and Local Ring Structure}
\label{p12-83-c1-s1:chap:app-hojas-algebra-local}

The \cref{p12-83-c1-s2:chap:app-completacion-gnomonica} constructed a single support
\(X_9\). APP has a single principal arithmetic matrix: nonadic multiplication
published by the positive digital root. Addition does not constitute a
second APP table; it is the internal accumulation law that generates the rows and
columns of the multiplicative matrix and provides an auxiliary observable for
comparison. This distinction will be maintained throughout the work.

\subsection{Nonadic multiplicative matrix and digital root}

For \(a,b\in\mathcal D_9=\{1,\ldots,9\}\), the APP entry is obtained
directly through
\begin{equation}
 \boxed{\Pi(a,b)=\operatorname{dr}_9(ab)
 =1+((ab-1)\bmod9).}
 \label{p12-83-c1-s1:eq:app-matriz-multiplicativa-directa}
\end{equation}
The only APP matrix is, therefore,
\[
\begin{array}{c|ccccccccc}
\Pi&1&2&3&4&5&6&7&8&9\\\hline
1&1&2&3&4&5&6&7&8&9\\
2&2&4&6&8&1&3&5&7&9\\
3&3&6&9&3&6&9&3&6&9\\
4&4&8&3&7&2&6&1&5&9\\
5&5&1&6&2&7&3&8&4&9\\
6&6&3&9&6&3&9&6&3&9\\
7&7&5&3&1&8&6&4&2&9\\
8&8&7&6&5&4&3&2&1&9\\
9&9&9&9&9&9&9&9&9&9
\end{array}
\]
Each row can also be reconstructed by reiterated addition of its index,
but that recurrence describes the internal law of the matrix and not a second
sheet. The rows and columns of indices \(1,2,4,5,7,8\) traverse the nine
classes; those of indices \(3,6,9\) are exactly those that contract the support.

In the residual chart, the set of units is
\[
 R_9^\times=(\mathbb Z/9\mathbb Z)^\times
   =\{[1],[2],[4],[5],[7],[8]\}.
\]
The maximal ideal is
\[
 \mathfrak m=(3)=\{[0],[3],[6]\}.
\]

\begin{theorem}[Classification of multiplicative rows]
\label{p12-83-c1-s1:thm:app-clasificacion-filas-producto}
For \(a\in\mathcal D_9\), the map
\[
 m_a:\mathcal D_9\longrightarrow\mathcal D_9,
 \qquad b\longmapsto\Pi(a,b)
\]
is bijective if and only if \([a]_9\in R_9^\times\). If
\([a]_9\in\mathfrak m\setminus\{0\}\), its image has three elements. If
\([a]_9=[0]_9\), its image has a single element.
\end{theorem}

\begin{proof}
The residual reduction of \(m_a\) is multiplication by \([a]_9\). In a
finite ring, that multiplication is bijective exactly when \([a]_9\)
is a unit. If \([a]_9=[3]\) or \([6]\), its image is the ideal
\(\mathfrak m\), of cardinality three. If \([a]_9=[0]\), every entry is sent
to the null class, published as \(9\).
\end{proof}

\begin{proposition}[Internality of the additive law]
\label{p12-83-c1-s1:prop:app-latinidad-aditiva}
The auxiliary observable
\[
 \Sigma(a,b)=\operatorname{dr}_9(a+b)
\]
acts by translations: for each \(a,c\in\mathcal D_9\) there exists a unique
\(b\in\mathcal D_9\) such that \(\Sigma(a,b)=c\), and analogously in the other
coordinate.
\end{proposition}

\begin{proof}
In the residual chart, \([a+b]_9=[c]_9\) has the unique solution
\([b]_9=[c-a]_9\). The positive section transports that class to a unique
element of \(\mathcal D_9\).
\end{proof}

The difference of geometry of fibers remains thus typed without duplicar the matrix
APP: the law additive internal only translates, whereas the matrix multiplicative
permutes on the units and contracts on the radical.

\begin{proposition}[Linear structure of the radical]
\label{p12-83-c1-s1:prop:app-radical-lineal}
The map
\[
 M_3:R_9\longrightarrow R_9,
 \qquad x\longmapsto3x
\]
satisfies
\[
 \ker M_3=\operatorname{im}M_3=\mathfrak m.
\]
The induced map
\[
 R_9/\mathfrak m\longrightarrow\mathfrak m,
 \qquad [x]\longmapsto3x
\]
is an isomorphism of vector spaces over \(\mathbb F_3\), but not a
ring isomorphism.
\end{proposition}

\begin{proof}
The congruence \(3x\equiv0\pmod9\) is equivalent to \(x\equiv0\pmod3\), whence
the kernel \(\mathfrak m\). The image is
\(\{[0],[3],[6]\}=\mathfrak m\). The first isomorphism theorem gives the
linear isomorphism. It is not a ring isomorphism because
\(\mathfrak m^2=0\) in \(R_9\), whereas the quotient
\(R_9/\mathfrak m\cong\mathbb F_3\) has a nonzero unit.
\end{proof}

The positive section publishes the non-trivial radical cycle as
\[
 369\longmapsto693\longmapsto936\longmapsto369,
\]
according to the position from which the triple is read. This notation preserves the
order of the coordinates; it does not turn the ideal into three independent
integers.

\subsection{Lifts and carry cocycle}

Let \(s_0:R_9\to\{0,\ldots,8\}\) be the normalized residual section. The
binary carry of the sum is
\[
 \kappa_0(a,b)
 =\frac{s_0(a)+s_0(b)-s_0(a+b)}9.
\]

\begin{lemma}[Cocyclic Identity]
\label{p12-83-c1-s1:lem:app-identidad-cociclo-acarreo}
For \(a,b,c\in R_9\),
\[
 \kappa_0(a,b)+\kappa_0(a+b,c)
 =\kappa_0(b,c)+\kappa_0(a,b+c).
\]
\end{lemma}

\begin{proof}
Multiply by nine and expand either of the two sides produces
\[
 s_0(a)+s_0(b)+s_0(c)-s_0(a+b+c).
\]
Therefore, both are equal.
\end{proof}

Define \(\chi_0(a)=1\) if \(a=[0]_9\) and \(0\) otherwise. Since
\[
 s_+(a)=s_0(a)+9\chi_0(a),
\]
the carryovers of the two sections are related by
\[
 \kappa_+(a,b)=
 \kappa_0(a,b)+\chi_0(a)+\chi_0(b)-\chi_0(a+b).
\]
Changing the representative modifies the cocycle by a coboundary. This is the
exact form for transporting memory when the null class is published as
\(9\) instead of \(0\).

\subsection{Integrated censuses of residues and quotients}

Summing over all pairs \((i,j)\in\mathcal D_9^2\), one obtains
\[
 \sum_{i,j}(i+j)=810,
 \qquad
 \sum_{i,j}ij=2025.
\]
The positive residues satisfy
\[
 \sum_{i,j}R^+_{i,j}=405,
 \qquad
 \sum_{i,j}R^\times_{i,j}=459.
\]
By the raised identities,
\[
 \sum_{i,j}Q^+_{i,j}=\frac{810-405}{9}=45,
 \qquad
 \sum_{i,j}Q^\times_{i,j}=\frac{2025-459}{9}=174.
\]

\begin{theorem}[Exact separation of the sum-product excess]
\label{p12-83-c1-s1:thm:app-separacion-exceso}
The integrated difference between product and sum decomposes uniquely
into a published part and a quotient part:
\[
 \boxed{
  2025-810
  =(459-405)+9(174-45)
  =54+9\cdot129.}
\]
\end{theorem}

\begin{proof}
We subtract, term to term, the identities
\(ij=R^\times_{i,j}+9Q^\times_{i,j}\) and
\(i+j=R^+_{i,j}+9Q^+_{i,j}\), and sum on the eighty and a pairs. The
censuses preceding dan the equality. The uniqueness proceeds of the division
Euclidean in each pair.
\end{proof}

The value \(54\) is not the complete difference between two supposed matrices.
It is the visible defect between the multiplicative matrix and the additive accumulator
in the positive chart. The term \(9\cdot129\) is the contribution
that remains in the quotients. Omitting it would confuse publication with
genealogy.

\subsection{Central block of \(2\) positions}

For \(k\in\mathbb Z/9\mathbb Z\), consideramos the block symmetric
\[
 B_k=
 \begin{pmatrix}
  k^2&k(k+1)\\
  k(k+1)&(k+1)^2
 \end{pmatrix}\pmod9.
\]
The diagonals coincide exactly when
\[
 k^2\equiv(k+1)^2\pmod9
 \quad\Longleftrightarrow\quad
 2k+1\equiv0\pmod9.
\]
Since \(2^{-1}=5\pmod9\), the unique solution is \(k=4\). Thus appears the
block
\[
 B_c=
 \begin{pmatrix}7&2\\2&7\end{pmatrix}.
\]
Let
\[
 R=\begin{pmatrix}0&1\\1&0\end{pmatrix},
 \qquad
 P_\pm=\frac12(I\pm R).
\]
Then
\[
 B_c=7I+2R=9P_++5P_-.
\]

\begin{theorem}[Uniqueness and spectrum of the central block]
\label{p12-83-c1-s1:thm:app-bloque-central}
The block \(B_c\) is the unique adjacent block in the family \(B_k\) with
equal diagonals. Its eigenvalues on the symmetric and
antisymmetric subspaces are, respectively, \(9\) and \(5\).
\end{theorem}

\begin{proof}
The congruence before proves the uniqueness. Since that \(R\) acts by
\(+1\) on \(\operatorname{im}P_+\) and by \(-1\) on
\(\operatorname{im}P_-\), the matrix \(7I+2R\) acts by \(9\) and \(5\),
respectively.
\end{proof}

The two ordered readings of its digits give
\[
 72+27=99,
 \qquad
 72-27=45=\det B_c,
\]
and the two diagonals give
\[
 77+22=99,
 \qquad
 77-22=55=54+1.
\]
These identities are internal controls of the block; they do not by themselves define the global defect.

\subsection{Mixed curvature of the \(2\) laws}

On the residual chart we write
\[
 S(x,y)=x+y,
 \qquad
 P(x,y)=xy.
\]
For a function \(T:X_9\to R_9\), let
\[
 \Delta_xT(x,y)=T(x+1,y)-T(x,y),
 \qquad
 \Delta_yT(x,y)=T(x,y+1)-T(x,y).
\]
We define the ordered curvature of a pair of laws by
\[
 \mathcal F(T_x,T_y)
 :=\Delta_x\Delta_yT_y-\Delta_y\Delta_xT_x.
\]

\begin{proposition}[Orientations of the mixed composition]
\label{p12-83-c1-s1:prop:app-curvatura-mixta}
The following hold
\[
 \mathcal F(S,S)=0,
 \qquad
 \mathcal F(P,P)=0,
\]
and, for mixed compositions,
\[
 \mathcal F(S,P)=+1,
 \qquad
 \mathcal F(P,S)=-1
 \quad\text{in }R_9.
\]
\end{proposition}

\begin{proof}
The mixed second differences of the linear law \(S\) vanish. For
\(P(x,y)=xy\), both pure mixed differences equal one and cancel when
compared in the same order. Combining a linear law and a bilinear one
leaves a single unit whose sign changes when the order of the sheets is reversed.
\end{proof}

The curvature not creates another matrix APP. Is recorded that the order
\(\Sigma\to\Pi\) and the order \(\Pi\to\Sigma\) are orientations opposite,
whereas the sectors pure remain flat. This triad
\(-1,0,+1\) will be typed in the unit TRIT.

\subsection{Obstructions of the arithmetic matrix of APP: translations, local radical, central block, mixed curvature, and conservation of carry}

\begin{criterion}[Refutation criteria for the APP matrix]
The unit is refuted if:
\begin{enumerate}
 \item the internal additive law ceases to act by translations;
 \item a nonunitary multiplicative row is declared bijective;
 \item \(R_9/\mathfrak m\) is identified with \(\mathfrak m\) as a ring;
 \item the exceso \(1215\) it replaces by \(54\) without conservar quotients;
 \item the central block is chosen without resolving \(2k+1=0\pmod9\);
 \item reversal of the mixed order does not change the sign of the curvature;
 \item the carry it eliminates to the cambiar between the sections \(s_0\) and
       \(s_+\).
\end{enumerate}
\end{criterion}

APP is now constructed as support, category of trajectories, nonadic
multiplicative matrix, internal additive law, local radical, carry cocycle,
and mixed orientation. The following unit
will not summarize these structures: it will take the triad of orientations
that has just emerged and construct the complete TRIT state with its three
quadratic regimes.


% Unidad U002. Redacción autónoma a partir del generador integral 1063/live.

\section{Gnomonic completion and arithmetic support of APP}
\label{p12-83-c1-s2:chap:app-completacion-gnomonica}

This unit constructs the APP support from the
nonadic cell of the \cref{p12-83-c1-s4:chap:segmentos-marcas-intersticios}. It does not presuppose a
continuous plane, an infinite lattice, or a numerical constant. The starting
object is finite: two ordered copies of the positive chart of
\(\mathbb Z/9\mathbb Z\). Completion of its interior interstices
produces exactly eighty-one states, and the residual translation law
turns that square support into a discrete torus.

Throughout the unit, APP denotes a mathematical structure and not an image.
Henceforth, the acronym APP names this structure; the properties that
follow depend on its objects and morphisms, not on a verbal expansion of the
acronym.

\subsection{Correspondence between \(2\) charts of the same residue class}

We define
\[
 \mathcal D_9:=\{1,2,\ldots,9\},
 \qquad
 \mathcal D_8:=\{1,2,\ldots,8\}.
\]
The residual publication application is
\[
 \lambda_9:\mathcal D_9\longrightarrow\mathbb Z/9\mathbb Z,
 \qquad
 \lambda_9(a)=[a]_9.
\]
Therefore, \(\lambda_9(9)=[0]_9\). Its inverse as a set map is the positive section
\[
 s_+([a]_9)=
 \begin{cases}
  a,&[a]_9\neq[0]_9,\quad 1\leq a\leq8,\\
  9,&[a]_9=[0]_9.
 \end{cases}
\]
The two expressions \(9\) and \([0]_9\) are not two phases: they are two
representatives of the same element in different charts.

For two coordinates we will write
\[
 \lambda_9^{\times2}:\mathcal D_9^2
       \longrightarrow X_9:=(\mathbb Z/9\mathbb Z)^2,
 \qquad
 (a,b)\longmapsto([a]_9,[b]_9).
\]
This map is a bijection of sets. The chart
\(\mathcal D_9^2\) makes the positive marks visible; the group \(X_9\)
makes the residual sum and the translations visible.

\begin{lemma}[State-preserving chart change]
\label{p12-83-c1-s2:lem:app-cambio-carta-sin-perdida}
The maps \(\lambda_9^{\times2}\) and
\(s_+^{\times2}\) are inverses. In particular, the positive representation and the
residual representation contain exactly the same eighty-one states.
\end{lemma}

\begin{proof}
In each coordinate,
\(\lambda_9\circ s_+=\operatorname{id}_{\mathbb Z/9\mathbb Z}\) and
\(s_+\circ\lambda_9=\operatorname{id}_{\mathcal D_9}\) hold. The Cartesian product of these identities gives the two required identities.
\end{proof}

\subsection{Interior kernel and gnomonic complement}

The inner core is
\[
 \mathcal N_8:=\mathcal D_8\times\mathcal D_8,
 \qquad |\mathcal N_8|=8^2=64.
\]
In the plane of indices, \(\mathcal N_8\) is the square that not uses the
mark of closure in no coordinate. For publish also the null class
it must completar the ninth row and the ninth column. We define the
complement gnomonic disjoint by
\[
 \mathcal G_9:=
  (\{9\}\times\mathcal D_9)
  \sqcup
  (\mathcal D_8\times\{9\}).
\]
The corner \((9,9)\) belongs to the first piece; for that reason the second uses
\(\mathcal D_8\) and not \(\mathcal D_9\). Consequently,
\[
 |\mathcal G_9|=9+8=17,
 \qquad
 \mathcal D_9^2=\mathcal N_8\sqcup\mathcal G_9,
 \qquad
 81=64+17.
\]

\begin{proposition}[Minimality of the Completion]
\label{p12-83-c1-s2:prop:app-minimalidad-completacion}
Let \(Y\) be a subset of \(\mathcal D_9^2\) that contains \(\mathcal N_8\) and whose projection onto each coordinate is all of \(\mathcal D_9\). If, in addition, \(Y\) contains all pairs required to translate one coordinate while the other remains arbitrary, then \(\mathcal G_9\subseteq Y\). Therefore, the addition of seventeen states is the minimal completion compatible with the two translations.
\end{proposition}

\begin{proof}
To transport the first coordinate through the closure class while
the second remains at an arbitrary value \(b\in\mathcal D_9\), all
pairs \((9,b)\) are necessary. Nine states are obtained. To
transport the second coordinate through closure while the first assumes
an arbitrary value \(a\in\mathcal D_8\), the remaining eight pairs
\((a,9)\) are necessary. The corner \((9,9)\) was already counted in the first
family. In total, \(9+8=17\) states are necessary, exactly those of
\(\mathcal G_9\).
\end{proof}

The hypothesis of compatibility with the two translations is essential. It would suffice
to add a single mark per coordinate for both projections to be
surjective, but that set would not contain the boundary states
needed to move one coordinate while holding the other fixed. APP requires the
complete product, not two disconnected axes.

\subsection{Exact realization of the interstitial interval via APP cells and boundary coordinates}

Let \(P_9\) be the ordered path with nine marks and eight interior edges
\[
 E(P_9)=\{e_1,\ldots,e_8\}.
\]
Closing the endpoints incorporates an edge \(e_9\) and produces the cycle
\(C_9\):
\[
 E(C_9)=E(P_9)\sqcup\{e_9\}.
\]
The pairs of interior interstices form \(E(P_9)^2\). The pairs of
interstices of the complete cycle form \(E(C_9)^2\). The difference is
\begin{align}
 E(C_9)^2\setminus E(P_9)^2
   &=(\{e_9\}\times E(C_9))
     \sqcup(E(P_9)\times\{e_9\}).
\label{p12-83-c1-s2:eq:app-complemento-intersticial}
\end{align}
The decomposition is disjoint because the corner \((e_9,e_9)\) is reserved for
the first piece.

\begin{theorem}[Equivalence of the mark model and the interstice model]
\label{p12-83-c1-s2:thm:app-equivalencia-modelos}
The map
\[
 \beta:E(C_9)^2\longrightarrow\mathcal D_9^2,
 \qquad
 \beta(e_i,e_j)=(i,j),
\]
is bijective. Its restrictions satisfy
\[
 \beta(E(P_9)^2)=\mathcal N_8,
 \qquad
 \beta(E(C_9)^2\setminus E(P_9)^2)=\mathcal G_9.
\]
Therefore, the identity \(81=64+17\) simultaneously counts pairs of
published marks and pairs of closed interstices.
\end{theorem}

\begin{proof}
Each edge of \(C_9\) has a unique index in \(\mathcal D_9\). The map \((i,j)\mapsto(e_i,e_j)\) is the explicit inverse of \(\beta\). If \(i,j\leq8\), the pair belongs to \(E(P_9)^2\) and its image belongs to \(\mathcal D_8^2\). If at least one of the indices is nine, the pair belongs to the complement of \cref{p12-83-c1-s2:eq:app-complemento-intersticial}; separating the case \(i=9\) from the case \(i\leq8,j=9\) produces exactly the two pieces of \(\mathcal G_9\).
\end{proof}

\subsection{Toroidal closure and translations}

Let
\[
 \mathbf e_1=([1]_9,[0]_9),
 \qquad
 \mathbf e_2=([0]_9,[1]_9).
\]
The set of elementary directions is
\[
 \mathscr D_{\rm APP}:=
 \{+\mathbf e_1,-\mathbf e_1,+\mathbf e_2,-\mathbf e_2\}.
\]
We define the oriented Cayley graph.
\[
 \Gamma_{\rm APP}:=\operatorname{Cay}(X_9,\mathscr D_{\rm APP}).
\]
Its vertices are the eighty-one elements of \(X_9\). For every
\(x\in X_9\) and \(d\in\mathscr D_{\rm APP}\), there is an oriented edge
\(x\xrightarrow{d}x+d\). The inverse edge bears the label \(-d\).

In the positive chart, unit translation is written by means of the published
successor and its carry:
\[
 u_+(a)=
 \begin{cases}
  a+1,&a<9,\\
  1,&a=9,
 \end{cases}
 \qquad
 \kappa_+(a)=
 \begin{cases}
  0,&a<9,\\
  1,&a=9.
 \end{cases}
\]
The raised identity is
\begin{equation}
 a+1=u_+(a)+9\kappa_+(a).
\label{p12-83-c1-s2:eq:app-sucesor-publicado-acarreo}
\end{equation}
Upon applying \(\lambda_9\), the carry term vanishes, leaving
\[
 \lambda_9(u_+(a))=\lambda_9(a)+[1]_9.
\]

\begin{lemma}[Compatibility between closure, residue, and elevation]
\label{p12-83-c1-s2:lem:app-compatibilidad-cierre-residuo}
The addition of \(e_9\), the identification \(9\sim[0]_9\), and the translation
\(x\mapsto x+\mathbf e_i\) describe three compatible operations in
distinct domains. The residual quotient publishes the return; the carry of
\cref{p12-83-c1-s2:eq:app-sucesor-publicado-acarreo} preserves the elevation that the
quotient omits.
\end{lemma}

\begin{proof}
Topological closure incorporates the edge joining the endpoints of the path. The
residual chart identifies the published endpoint with the null class. If
\(a<9\), the successor has zero carry; if \(a=9\), the lifted equality
is \(10=1+9\). Both cases descend to the same residual sum, but the
value of \(\kappa_+\) distinguishes them.
\end{proof}

\begin{corollary}[Toroidal support]
\label{p12-83-c1-s2:cor:app-soporte-toroidal}
Identifying opposite sides of the chart \(\mathcal D_9^2\) yields a
geometric realization of the group \(X_9\). Each elementary translation is
defined at every vertex and has an inverse.
\end{corollary}

\subsection{APP Path Category: Objects, Prolongation Morphisms and Compatible Composition}

\begin{definition}[Oriented trajectory]
A trajectory of length \(n\geq0\) is a pair
\[
 \gamma=(x_0;d_1,\ldots,d_n),
 \qquad
 x_0\in X_9,\qquad d_k\in\mathscr D_{\rm APP}.
\]
Its vertices are
\[
 x_k=x_0+\sum_{r=1}^k d_r,
 \qquad 0\leq k\leq n.
\]
The source is \(s(\gamma)=x_0\) and the target is
\(t(\gamma)=x_n\). For \(n=0\) the identity trajectory in
\(x_0\) is obtained.
\end{definition}

If \(t(\gamma)=s(\delta)\), its concatenation is
\[
 \delta\circ\gamma
  =(x_0;d_1,\ldots,d_n,e_1,\ldots,e_m),
\]
where \(e_1,\ldots,e_m\) are the directions of \(\delta\). Concatenation is associative, and paths of zero length are identities. We thus obtain the free category
\[
 \mathsf{Path}(\Gamma_{\rm APP}).
\]

\begin{proposition}[Finiteness of Support and Countability of Trajectories]
\label{p12-83-c1-s2:prop:app-soporte-finito-trayectorias-numerables}
The APP state set is finite, of cardinality \(81\). For a fixed source,
there are \(4^n\) direction words of length \(n\). The union
of all finite trajectories is countable.
\end{proposition}

\begin{proof}
The first assertion follows from \(|X_9|=9^2\). At each of the \(n\)
steps there are four independent choices, whence \(4^n\). Finally,
\(\bigcup_{n\geq0}X_9\times\mathscr D_{\rm APP}^n\) is a countable union
of finite sets.
\end{proof}

This proposition separates two levels that will not be identified: APP has a
finite support, but admits trajectories of arbitrary length. The
prolongation of a route does not create new vertices of the support; it creates new
genealogy over them.

\subsection{Multiplicative matrix and internal additive law on a single support}

For \(n\in\mathbb Z_{>0}\), the root digital positive is
\[
 \operatorname{dr}_9(n):=1+((n-1)\bmod9).
\]
On the positive chart we define the principal matrix and the internal law.
\[
 \Sigma(a,b):=\operatorname{dr}_9(a+b),
 \qquad
 \Pi(a,b):=\operatorname{dr}_9(ab).
\]
Both maps have the same domain \(\mathcal D_9^2\) and the same
codomain \(\mathcal D_9\), but they do not have the same status: \(\Pi\) publishes
the unique APP matrix, and \(\Sigma\) generates its accumulations. Their residual versions
are
\[
 \bar\Sigma([a]_9,[b]_9)=[a+b]_9,
 \qquad
 \bar\Pi([a]_9,[b]_9)=[ab]_9.
\]
The change of chart is compatible with both operations:
\begin{align*}
 \lambda_9\circ\Sigma
   &=\bar\Sigma\circ\lambda_9^{\times2},\\
 \lambda_9\circ\Pi
   &=\bar\Pi\circ\lambda_9^{\times2}.
\end{align*}

To preserve the integer quotient, lifts are introduced.
\begin{align}
 a+b &= R^+_{a,b}+9Q^+_{a,b},
 &1\leq R^+_{a,b}\leq9,
\label{p12-83-c1-s2:eq:app-lift-suma}\\
 ab &= R^\times_{a,b}+9Q^\times_{a,b},
 &1\leq R^\times_{a,b}\leq9.
\label{p12-83-c1-s2:eq:app-lift-producto}
\end{align}
Here
\[
 R^+_{a,b}=\Sigma(a,b),
 \qquad
 R^\times_{a,b}=\Pi(a,b),
\]
and the quotients are uniquely determined by
\[
 Q^+_{a,b}=\frac{a+b-R^+_{a,b}}9,
 \qquad
 Q^\times_{a,b}=\frac{ab-R^\times_{a,b}}9.
\]

\begin{definition}[APP genealogical]
\label{p12-83-c1-s2:def:app-genealogica-autonoma}
The APP structure at this stage is the tuple
\[
 \operatorname{APP}:=
 \bigl(
  X_9,
  \Gamma_{\rm APP},
  \mathsf{Path}(\Gamma_{\rm APP}),
  \Sigma,\Pi,
  R^+,Q^+,R^\times,Q^\times
 \bigr).
\]
The genealogy of a trajectory preserves, for each vertex visited, the
position, incoming direction, multiplicative entry, additive accumulation,
and their respective residues and quotients.
\end{definition}

\begin{theorem}[Causal construction of the APP support]
\label{p12-83-c1-s2:thm:app-cadena-constructiva-autonoma}
The chain
\[
 \mathcal D_9
 \longrightarrow
 \mathcal D_8^2
 \longrightarrow
 \mathcal D_8^2\sqcup\mathcal G_9
 \longrightarrow
 \mathcal D_9^2
 \xrightarrow{\lambda_9^{\times2}}
 X_9
 \longrightarrow
 \Gamma_{\rm APP}
 \longrightarrow
 \mathsf{Path}(\Gamma_{\rm APP})
 \longrightarrow
 (\Sigma,\Pi)
\]
constructs, without loss of cardinality, the support, motion, paths, APP multiplicative matrix and its internal additive law.
\end{theorem}

\begin{proof}
The \cref{p12-83-c1-s2:prop:app-minimalidad-completacion} constructs the complete chart from
the kernel. The \cref{p12-83-c1-s2:thm:app-equivalencia-modelos} identifies its
interstitial model. The \cref{p12-83-c1-s2:lem:app-cambio-carta-sin-perdida} transports
the eighty-one states to \(X_9\). The Cayley graph incorporates the four
reversible translations without changing the vertices. The free category
incorporates all finite concatenations. Finally, \(\Sigma\) and
\(\Pi\) are distinct total maps on the same domain; the
equations \cref{p12-83-c1-s2:eq:app-lift-suma,p12-83-c1-s2:eq:app-lift-producto} preserve the part
that residual reduction does not publish.
\end{proof}

\subsection{Decomposition of the APP orthogram into kernel \(8\times8\), gnomon of \(17\) states, and chart \(9\times9\) with opposite sides identified}

The \cref{fig:app-ortograma} must be read in three layers: the inner square
of sixty-four states, the gnomon of seventeen states, and the
identification of opposite sides. The nine-by-nine chart thereby completes the
eight-by-eight core by means of the closure gnomon. The multiplicative
matrix and the internal additive law are evaluated at every vertex of the same
state, preserve distinct residue--quotient lifts, and do not constitute
two APP squares.

\subsection{Obstructions to the fidelity of the APP support: cardinals \(81/17\), reversible translations, and separation of the sum and product sheets}

\begin{criterion}[Refutation criteria for the APP support]
The construction is refuted if any of the following facts
occurs:
\begin{enumerate}
 \item the complete support has cardinality other than \(81\);
 \item the gnomonic complement does not have cardinal \(17\);
 \item a state of \(X_9\) lacks any of the four translations;
 \item it identify \(\Sigma\) and \(\Pi\) because coincide in a value;
 \item \(Q^+\) or \(Q^\times\) is eliminated upon returning through a boundary;
 \item trajectories of different lengths are presented as new
       support vertices;
 \item it interpreta \(9\) and \([0]_9\) as phases independent.
\end{enumerate}
\end{criterion}

The unit has constructed APP up to its complete genealogical evaluation. It has not yet
oriented the difference between the matrix and its accumulator, has not
introduced the TRIT,
and has not defined the Topological Phase Kernel. The following unit will study
the algebraic structure of rows and columns, the sector of units, and the
carry cocycle that prepares that orientation.


% Unidad U005. Redacción autónoma desde 03c/03h/20 de la autoridad 1063/live.

\long\def\HMTDeferredTPKBase{%
\section{Topological Phase Kernel: observable base and enriched state}
\label{p12-83-c1-s3:chap:tpk-categoria-estado-enriquecido}

APP has supplied a support, a multiplicative matrix, and an internal
additive law; TRIT has
supplied a lifted orientation. The Topological Phase Kernel, hereinafter
TPK, specifies how those data are transported jointly, which
coordinates are preserved when a route is closed, and which readers can be published
without identifying the publication with the state. Its mathematical object is a
category over a base of paths, endowed with an explicit family of
fibers and prolongation relations.

\subsection{The double base of APP paths}

Let \(\mathsf P_{\rm APP}\) be the category libre constructed in the
\cref{p12-83-c1-s2:chap:app-completacion-gnomonica}. We define
\[
 \mathsf P_{\rm APP}^{(2)}
 :=\mathsf P_{\rm APP}\times\mathsf P_{\rm APP}.
\]
An object is a pair of positions
\((p^\Sigma,p^\Pi)\in X_9^2\). A morphism is a pair of trajectories; one
of them may be the identity. The first component is evaluated by
\(\Sigma\), the second by \(\Pi\). They are not two lattices: they are two typed
cursors over the same APP support.

The set of cardinal directions is
\[
 \mathscr D_{\rm card}:=\{N,E,S,O\},
\]
with shifts
\[
 \delta(N)=(-1,0),\quad
 \delta(E)=(0,1),\quad
 \delta(S)=(1,0),\quad
 \delta(O)=(0,-1)
 \quad\text{in }X_9.
\]
The orientation involution is
\[
 \overline N=S,
 \quad \overline S=N,
 \quad \overline E=O,
 \quad \overline O=E.
\]

\subsection{Nonadic phase, dodecaphasic sector and calendar}

The complete clock of this stage is
\[
 \Theta_{108}:=\mathbb Z/108\mathbb Z.
\]
For \(\theta\in\Theta_{108}\), let
\(\widetilde\theta\in\{0,\ldots,107\}\) be its normalized representative.
We define three distinct coordinates:
\begin{align*}
 g_0(\theta)&=[\widetilde\theta+1]_9\in C_9,\\
 r(\theta)&=1+(\widetilde\theta\bmod9)\in\mathcal D_9,\\
 \beta(\theta)&=
 \left[\left\lfloor\frac{\widetilde\theta}{9}\right\rfloor\right]_{12}
 \in C_{12}.
\end{align*}
The first is the phase class, the second its positive label, and the third
the dodecaphasic sector. The map \(\beta\) is a block map;
it is not the canonical reduction \(\theta\bmod12\).

The signature and the active sheet are
\[
 \tau(\theta):=\varepsilon_9(r(\theta))\in\mathsf{TRIT},
\]
and
\[
 \mu(\theta)=
 \begin{cases}
  \Sigma,&\tau(\theta)=+1,\\
  \Pi,&\tau(\theta)=-1,\\
  \varnothing,&\tau(\theta)=0.
 \end{cases}
\]
The symbol \(\varnothing\) denotes a threshold step in which neither cursor
moves; it does not denote the absence of a state.

\begin{definition}[Observable configuration]
\label{p12-83-c1-s3:def:tpk-configuracion-observable}
The observable base is
\[
 \mathcal Q_{\rm obs}
 :=X_9^2\times\mathscr D_{\rm card}^2\times\Theta_{108}.
\]
An element is written
\[
 q=(p^\Sigma,p^\Pi,d^\Sigma,d^\Pi,\theta).
\]
It contains two positions, two cardinal orientations, and a calendar index.
\end{definition}

\subsection{Sheet selector \(M_{\rm ph}:\mathcal D_9\to\mathsf{TRIT}\) and joint evolution of position, orientation, and calendar}

The function
\[
 M_{\rm ph}:\mathcal D_9\longrightarrow\mathsf{TRIT},
 \qquad M_{\rm ph}(r)=\varepsilon_9(r)
\]
is the internal sheet selector. It selects which evaluation acts; it does not by itself
displace any position.

We define the orientation inversion indicator by
\[
 \eta_{27}(\theta)=
 \begin{cases}
  -1,&\widetilde\theta+1\in\{27,54,81,108\},\\
  +1,&\text{otherwise}.
 \end{cases}
\]
The new directions are
\[
 d^{\nu,+}=
 \begin{cases}
  \overline{d^\nu},&\eta_{27}(\theta)=-1,\\
  d^\nu,&\eta_{27}(\theta)=+1,
 \end{cases}
 \qquad \nu\in\{\Sigma,\Pi\}.
\]
The positions evolve via
\[
 (p^{\Sigma,+},p^{\Pi,+})=
 \begin{cases}
 (p^\Sigma+\delta(d^\Sigma),p^\Pi),&\tau(\theta)=+1,\\
 (p^\Sigma,p^\Pi+\delta(d^\Pi)),&\tau(\theta)=-1,\\
 (p^\Sigma,p^\Pi),&\tau(\theta)=0.
 \end{cases}
\]

\begin{definition}[Observable evolution]
\label{p12-83-c1-s3:def:tpk-evolucion-observable}
The map
\[
 F_{\rm obs}:\mathcal Q_{\rm obs}\longrightarrow\mathcal Q_{\rm obs}
\]
is defined by
\[
 F_{\rm obs}(q)=
 (p^{\Sigma,+},p^{\Pi,+},d^{\Sigma,+},d^{\Pi,+},[\theta+1]_{108}).
\]
The category libre of the graph \(q\to F_{\rm obs}(q)\) it is denoted by
\(\mathsf P_{\rm obs}\).
\end{definition}

Each arrow of \(\mathsf P_{\rm obs}\) induces a pair of trajectories in
APP, so there exists a functor
\[
 B:\mathsf P_{\rm obs}\longrightarrow
 \mathsf P_{\rm APP}^{(2)}.
\]
If \(\tau=+1\), the second component of \(B\) is the identity; if
\(\tau=-1\), the first is the identity; if \(\tau=0\), both are the identity.

\begin{proposition}[Separation between selection, motion, and remark]
\label{p12-83-c1-s3:prop:tpk-separacion-seleccion-movimiento}
The selector \(M_{\rm ph}\), the translations
\(T_d(p)=p+\delta(d)\), and any temporal sampling have distinct domains and
codomains. They are not interchangeable operators.
\end{proposition}

\begin{proof}
\(M_{\rm ph}\) takes a phase label and returns a tritic orientation.
\(T_d\) is an automorphism of \(X_9\). A sampling acts on a
sequence or a category of arrows. Identifying two of them would violate
typing. \(F_{\rm obs}\) composes them in a declared order without erasing that
distinction.
\end{proof}

\subsection{Decomposition of the enriched state in \(6\) families fibradas}

Over each \(q\in\mathcal Q_{\rm obs}\) we fix the total fiber
\[
 \mathcal F_q=
 \mathsf O_q\times
 \mathsf H_q\times
 \mathsf C_q\times
 \mathsf S_q\times
 \mathsf B_q\times
 \Lambda_q.
\]
Each factor has its own type and function.

\subsubsection{Joint determination of the orientation and sheet of the enriched state}

\[
 \mathsf O_q\subseteq
 \mathscr D_{\rm card}^2\times\{+1,-1\}.
\]
The two directions coincide with those stored in \(q\); the last
coordinate records the sheet orientation when the realization requires it.
The sheet label does not replace the cardinal directions.

\subsubsection{Euclidean decomposition of the state into compatible residue and quotient}

The first lift fiber is
\[
 \mathsf H^{(1)}=\mathcal D_9\times\mathbb Z,
 \qquad
 \operatorname{rec}_9(r,u)=r+9u.
\]
For each integer \(n\) there exist unique
\[
 r=\operatorname{dr}_9(n),
 \qquad
 u=\frac{n-r}{9}.
\]
A greater depth consists of a declared tower of these fibers and their
truncation maps. The term ``Hensel'' does not replace the definition of
that tower.

\subsubsection{Carry operator and joint update of residue, quotient, and memory}

\(\mathsf C_q\) is an abelian group. In the readers of decimal blocks
one takes \(\mathsf C_q=\mathbb Z\). The increment is not an independent
coordinate, but a cochain over the arrows:
\[
 \kappa(r)\in\mathsf C_{q'}
 \quad\text{for }r:q\to q'.
\]

\subsubsection{Construction of the full signature from orientation, carry, and boundary}

The signature is derived via a typed application.
\[
 \operatorname{Sig}_q:
 \mathsf H_q\times\mathsf C_q\times\Lambda_q
 \longrightarrow\mathsf S_q.
\]
In the nonadic boundary realization, individual conservation is maintained,
\[
 q_\partial,a_\partial,c_\partial\in\mathbb F_3^6,
 \qquad
 \operatorname{colw}\in\mathbb Z_{\geq0}^{6},
 \qquad
 \rho_W,r_\partial\in\mathbb F_3^3.
\]
These coordinates are margin, axial combination, cubic condition, column
weights, Witt residual class, and ternary return. They do not condense into a
single word if a subsequent proof uses any of them.

\subsubsection{Cylinder and boundary}

\[
 \mathsf B_q:=\mathsf{Cyl}_q\times\mathsf{Fr}_q.
\]
The cylinder localizes a family of prefixes; the boundary determines which
extensions are compatible. Survival will be a relation between
states endowed with both data, not a number added a posteriori.

\subsubsection{Genealogical transport record}

\(\Lambda_q\) contains the transport records whose endpoint is
\(q\). A prolongation \(r:q\to q'\) acts by concatenation:
\[
 \lambda\longmapsto r\circ\lambda.
\]
The genealogical register preserves the route; its observable endpoint does not determine it.

\subsection{Fiber transport}

For each arrow \(r:q\to q'\) maps are declared
\[
 \mathsf H(r):\mathsf H_q\to\mathsf H_{q'},
 \quad
 \mathsf C(r):\mathsf C_q\to\mathsf C_{q'},
 \quad
 \Lambda(r):\Lambda_q\to\Lambda_{q'}.
\]
They respect identities and composition, and \(\mathsf C(r)\) is a homomorphism.
The carry satisfies
\begin{equation}
 \kappa(r_2\circ r_1)
 =\kappa(r_2)+\mathsf C(r_2)\bigl(\kappa(r_1)\bigr).
\label{p12-83-c1-s3:eq:tpk-cocadena-transporte}
\end{equation}
Therefore, the transport of the determined coordinates is
\begin{align*}
 h'&=\mathsf H(r)(h),\\
 c'&=\mathsf C(r)(c)+\kappa(r),\\
 \lambda'&=\Lambda(r)(\lambda),\\
 \sigma'&=\operatorname{Sig}_{q'}(h',c',\lambda').
\end{align*}

\begin{lemma}[Associativity of the transported carry]
\label{p12-83-c1-s3:lem:tpk-asociatividad-acarreo}
The equation \cref{p12-83-c1-s3:eq:tpk-cocadena-transporte} guarantees that transporting a
state by \(r_1\) and then by \(r_2\) produces the same carry as
transporting it by \(r_2\circ r_1\).
\end{lemma}

\begin{proof}
The successive transport gives
\[
 \mathsf C(r_2)\bigl(\mathsf C(r_1)c+\kappa(r_1)\bigr)+\kappa(r_2).
\]
By functoriality of \(\mathsf C\), the first term is
\(\mathsf C(r_2r_1)c\); the sum of the other two is
\(\kappa(r_2r_1)\) by \cref{p12-83-c1-s3:eq:tpk-cocadena-transporte}.
\end{proof}

\subsection{Extension relations}

Orientation and boundary are not imposed by means of a universal function.
For each arrow \(r:q\to q'\), a relation is declared
\[
 \operatorname{Ext}_r\subseteq\mathcal F_q\times\mathcal F_{q'}.
\]
The two conditions are required.
\[
 \Delta_q\subseteq\operatorname{Ext}_{\operatorname{id}_q},
 \qquad
 \operatorname{Ext}_{r_2}\circ\operatorname{Ext}_{r_1}
 \subseteq\operatorname{Ext}_{r_2\circ r_1}.
\]
The first preserves identities; the second preserves composition. A
relation may admit several local extensions and leave only one after a
larger horizon.

\begin{definition}[The Topological Phase Kernel]
\label{p12-83-c1-s3:def:tpk-categoria-autonoma}
An object of the TPK is a tuple
\[
 x=(q,o,h,c,\sigma,C,b,\lambda)
\]
such that
\[
 q\in\mathcal Q_{\rm obs},\quad
 o\in\mathsf O_q,\quad
 h\in\mathsf H_q,\quad
 c\in\mathsf C_q,\quad
 \lambda\in\Lambda_q,
\]
\[
 \sigma=\operatorname{Sig}_q(h,c,\lambda),
 \qquad
 (C,b)\in\mathsf{Cyl}_q\times\mathsf{Fr}_q.
\]
We associate to each object its fiber component.
\[
 \mathbf f(x):=(o,h,c,\sigma,(C,b),\lambda)\in\mathcal F_q.
\]
A morphism \(x\to x'\) is a triple \((r,x,x')\) in which
\(r:q\to q'\) belongs to \(\mathsf P_{\rm obs}\), the transportable
coordinates satisfy the preceding laws, and
\((\mathbf f(x),\mathbf f(x'))\in
\operatorname{Ext}_r\). Composition is
\[
 (r_2,x',x'')\circ(r_1,x,x')=(r_2r_1,x,x''),
\]
and the identity of \(x\) is \((\operatorname{id}_q,x,x)\).
\end{definition}

\begin{theorem}[Categorical structure of the TPK]
\label{p12-83-c1-s3:thm:tpk-es-categoria}
The \cref{p12-83-c1-s3:def:tpk-categoria-autonoma} defines a category and the map
\[
 p:\mathsf{TPK}\longrightarrow\mathsf P_{\rm obs},
 \qquad p(x)=q,\qquad p(r,x,x')=r,
\]
is a functor.
\end{theorem}

\begin{proof}
Route composition is associative. The
\cref{p12-83-c1-s3:lem:tpk-asociatividad-acarreo} proves compatibility of the
carry coordinate; the declared functoriality of \(\mathsf H\),
\(\mathsf C\), and \(\Lambda\) proves compatibility of the other determined
coordinates. The stability of \(\operatorname{Ext}\) guarantees that the
composite remains admissible, and its diagonal condition provides the
identities. The map \(p\) preserves identities and composition by
construction.
\end{proof}

\subsection{Projections and typed losses}

The elementary projections are
\[
 \phi(x)=g_0(\theta),
 \qquad
 \tau(x)=\varepsilon_9(r(\theta)),
 \qquad
 \mu(x)=\mu(\theta).
\]
The projection to the local TRIT state is
\[
 \pi_{\rm TRIT}(x)
 =\bigl(\tau,\mu,o,h,c,\sigma,\lambda,\phi\bigr).
\]
It forgets the cylinder, the boundary, and part of \(q\); it does not define another container.
Two states may have the same image under \(\pi_{\rm TRIT}\) and differ
in the forgotten information.

\subsection{Criterion of route survival via conservation of the projected boundary}

Let \(\mathcal X_{\rm TPK}^{(n)}\) be the family of states at depth
\(n\). For \(N\geq n\), we define
\[
 \operatorname{Surv}_n^{(N)}=
 \left\{x_n:
 \begin{array}{l}
 \text{there exist }x_{n+1},\ldots,x_N\text{ with}\\
 (x_j,x_{j+1})\in\operatorname{Ext}_{r_j}
 \text{ for }n\leq j<N
 \end{array}\right\}.
\]

\begin{proposition}[Survival filtration]
\label{p12-83-c1-s3:prop:tpk-filtracion-supervivencia}
For \(N\geq n\),
\[
 \operatorname{Surv}_n^{(N+1)}
 \subseteq
 \operatorname{Surv}_n^{(N)}.
\]
If the scheme is defined to all depth, the states
indefinitely prolongable are
\[
 \operatorname{Surv}_n
 =\bigcap_{N\geq n}\operatorname{Surv}_n^{(N)}.
\]
\end{proposition}

\begin{proof}
A prolongation to \(N+1\) truncates to a prolongation to \(N\),
which proves the inclusion. The intersection expresses exactly
prolongability for every finite horizon.
\end{proof}

\subsection{Operator tree of the container}

\begin{figure}[htbp]
 \centering
 \begin{tikzpicture}[
  node distance=7mm and 9mm,
  every node/.style={font=\small,align=center},
  root/.style={draw=black,very thick,rounded corners,inner sep=5pt},
  op/.style={draw=black,rounded corners,inner sep=4pt},
  edge/.style={-{Latex[length=2mm]},semithick}
 ]
  \node[root] (tpk) {TPK\\$\mathcal X_{\rm TPK}^{\rm enr}$};
  \node[op,below left=of tpk] (sel) {selection\\$M_{\rm ph}$};
  \node[op,below=of tpk] (mov) {transport\\$T_d,F_{\rm obs}$};
  \node[op,below right=of tpk] (fib) {fibers\\$\mathsf H,\mathsf C,\mathsf S,
    \mathsf B,\Lambda$};
  \node[op,below=17mm of mov] (cal) {calendar\\$C_{12}\hookrightarrow
    C_{108}\twoheadrightarrow C_9$};
  \node[op,below=of cal] (ext) {extension, cylinders and survival};
  \draw[edge] (tpk) -- (sel);
  \draw[edge] (tpk) -- (mov);
  \draw[edge] (tpk) -- (fib);
  \draw[edge] (sel) |- (cal);
  \draw[edge] (mov) -- (cal);
  \draw[edge] (fib) |- (cal);
  \draw[edge] (cal) -- (ext);
 \end{tikzpicture}
 \caption{Operator containment of the Topological Phase Kernel. The selector,
 transport, readers, and fibers are not parallel containers: they are
 maps or coordinates of the same enriched state.}
 \label{p12-83-c1-s3:fig:tpk-arbol-operatorio-autonomo}
\end{figure}

\subsection{Obstructions of the enriched state of the TPK: typing of the selector, carry cochain, compatible extension, and genealogical conservation}

\begin{criterion}[Refutation Criteria for the Enriched State]
The definition is refuted if:
\begin{enumerate}
 \item the selector \(M_{\rm ph}\) it usa as translation;
 \item \(\theta\), \(g_0(\theta)\), \(r(\theta)\), and
       \(\beta(\theta)\) are identified;
 \item a signature is allowed to be chosen independently of
       \((h,c,\lambda)\);
 \item the composition of carries incumple
       \cref{p12-83-c1-s3:eq:tpk-cocadena-transporte};
 \item \(\operatorname{Ext}_r\) is replaced by an unproved functional choice;
 \item a projection that has forgotten the cylinder, boundary,
       or genealogical record is called a complete state;
 \item the phase return erases any of the coordinates of fiber.
\end{enumerate}
\end{criterion}

This unit has defined the container. The next will enumerate its readers,
operators, odometer, and calendars, and will prove why the returns of
nine, twenty-seven, and one hundred eight steps are not restarts of the state.
}%


% Unidad U001. Redacción autónoma; no procede del borrador condensado.

\section{Marked segments, interstices and the nonadic cell}
\label{p12-83-c1-s4:chap:segmentos-marcas-intersticios}

The construction begins before every constant and before every choice of
physical unit. The inaugural datum is not an already constituted real number, but an
ordered segment, a family of marks, and the adjacency relation between
consecutive marks. This separation is necessary because a mark, an
interval, the length of that interval, and the numeral that publishes that
length are objects of different types.

\subsection{Finite positional chart and decimal residual projection}

\begin{definition}[Marked segment of order ten]
Let
\[
 \mathsf M_{10}:=\{m_0,m_1,\ldots,m_{10}\}
\]
a set totally ordered by
\(m_0<m_1<\cdots<m_{10}\). Its adjacency graph is the path
\[
 \mathsf P_{10}:\quad
 m_0\longrightarrow m_1\longrightarrow\cdots\longrightarrow m_{10}.
\]
We shall denote by
\[
 \mathsf I_k:=(m_k,m_{k+1}),\qquad 0\leq k\leq9,
\]
the abstract interstice between two consecutive marks. The
family of interstices is
\(\mathsf E_{10}:=\{\mathsf I_0,\ldots,\mathsf I_9\}\).
\end{definition}

The open notation \((m_k,m_{k+1})\) does not yet presuppose the real line. It
indicates only the ordered pair of endpoints and the edge joining them.
When the Archimedean chart is chosen
\[
 \chi_{10}(m_k)=k,
\]
the abstract interstice \(\mathsf I_k\) is published by the semiopen interval
\([k,k+1)\). Therefore, the index \(k\), the edge
\(\mathsf I_k\), and the interval \([k,k+1)\) will not be identified in the text:
the first is a label, the second is a combinatorial element, and the
third is its realization in a chart.

\begin{lemma}[Census of Marks and Gaps]
\label{p12-83-c1-s4:lem:censo-marcas-intersticios}
The segment \(\mathsf P_{10}\) contains eleven marks and ten interstices. More generally, a path with \(n+1\) marks has exactly \(n\) edges.
\end{lemma}

\begin{proof}
The map
\[
 \{0,\ldots,n-1\}\longrightarrow E(\mathsf P_n),
 \qquad k\longmapsto(m_k,m_{k+1}),
\]
is bijective. The particular assertion is obtained with \(n=10\).
\end{proof}

The endpoint \(m_0\) fixes the origin of the chart, and \(m_{10}\) fixes its threshold. Neither by itself constitutes an additional interstice. This elementary remark prevents three subsequent confusions: counting endpoints as cells, confusing closure with a new phase state, and treating zero as a tenth phase independent of the nonadic cycle.

\subsection{Extraction of the cell of \(9\) phases}

The decimal chart distinguishes the anchor \([0,1)\) of the periodic cell. The
nine positive marks
\[
 \mathsf D_9:=\{1,2,3,4,5,6,7,8,9\}
\]
it utilizaran as representantes publicos of
\(\mathbb Z/9\mathbb Z\), haciendo corresponder the mark \(9\) with the class
residual \(\overline 0\). We define
\[
 \lambda_9:\mathsf D_9\longrightarrow\mathbb Z/9\mathbb Z,
 \qquad \lambda_9(a)=\overline a.
\]

\begin{proposition}[Positive chart of the nonadic classes]
\label{p12-83-c1-s4:prop:carta-positiva-nonadica}
The map \(\lambda_9\) is bijective. Its inverse is the positive section
\[
 s_+(\overline a)=
 \begin{cases}
 a,&1\leq a\leq8,\\
 9,&\overline a=\overline0.
 \end{cases}
\]
In particular, the public representation \(9\) and the residual representation \(0\)
represent the same class, but belong to different charts.
\end{proposition}

\begin{proof}
The nine classes of \(\mathbb Z/9\mathbb Z\) have unique representatives in
\(\{0,1,\ldots,8\}\). Replacing the representative \(0\) by \(9\) produces
another complete transversal, so \(\lambda_9\) is bijective and \(s_+\)
is its inverse.
\end{proof}

The phase cell is not obtained by arbitrarily erasing the origin. The segment
\([0,1)\) fixes the decimal chart from which the cycle is observed; the nine
phases are the nine classes published as \(1,\ldots,9\). In the lifted
chart, the step that closes the cell is represented by \([9,10)\); in the
residual chart, that same step returns from \(\overline0\) to \(\overline1\).
Closure and change of representative are not the same operation.

\subsection{Position, height and Euclidean division}

\begin{definition}[Nonadic Coordinates]
For each integer \(n\geq1\), we define
\[
 \rho_9(n):=1+((n-1)\bmod9),
 \qquad
 q_9(n):=\left\lfloor\frac{n-1}{9}\right\rfloor.
\]
The first coordinate is called phase position; the second, turn height or memory quotient.
\end{definition}

\begin{theorem}[Unique nonadic decomposition]
\label{p12-83-c1-s4:thm:descomposicion-nonadica-unica}
Every integer \(n\geq1\) admits a unique expression
\[
 n=9q+\rho,
 \qquad q\in\mathbb Z_{\geq0},\quad \rho\in\mathsf D_9.
\]
In esa escritura \(q=q_9(n)\) and \(\rho=\rho_9(n)\).
\end{theorem}

\begin{proof}
We apply Euclidean division to \(n-1\): there exist unique
\(q\geq0\) and \(r\in\{0,\ldots,8\}\) such that
\(n-1=9q+r\). Taking \(\rho=r+1\) gives
\(n=9q+\rho\), with \(1\leq\rho\leq9\). The uniqueness of \(q,r\) implies that
of \(q,\rho\).
\end{proof}

\begin{corollary}[Phase return with memory increment]
\label{p12-83-c1-s4:cor:retorno-fase-incremento-memoria}
For every \(n\geq1\),
\[
 \rho_9(n+9)=\rho_9(n),
 \qquad
 q_9(n+9)=q_9(n)+1.
\]
Consequently, a full rotation preserves the visible phase and changes the raised state.
\end{corollary}

\begin{proof}
If \(n=9q+\rho\), then \(n+9=9(q+1)+\rho\); uniqueness of
\cref{p12-83-c1-s4:thm:descomposicion-nonadica-unica} gives the two identities.
\end{proof}

The one-step evolution is written unambiguously as
\[
 (q,\rho)\longmapsto
 \begin{cases}
 (q,\rho+1),&1\leq\rho<9,\\
 (q+1,1),&\rho=9.
 \end{cases}
\]
The second line is not a restart. It recovers the initial phase label and
increases the quotient that records how many turns have been performed.

\subsection{Circular projection periodicity and injectivity of the nonadic helical lifting}

Let
\[
 \vartheta_9(\rho)=\frac{2\pi(\rho-1)}9.
\]
The phase projection onto the unit circle is
\[
 \mathcal R_9(n)=
 \bigl(\cos\vartheta_9(\rho_9(n)),
       \sin\vartheta_9(\rho_9(n))\bigr).
\]
To simultaneously preserve position and height we introduce the helical immersion.
\[
 \mathcal H_9(n)=
 \left(
  \cos\vartheta_9(\rho_9(n)),
  \sin\vartheta_9(\rho_9(n)),
  q_9(n)+\frac{\rho_9(n)-1}{9}
 \right).
\]

\begin{proposition}[The circular projection does not determine the state]
\label{p12-83-c1-s4:prop:rueda-no-determina-estado}
The map \(\mathcal R_9\) is periodic with period nine. The map
\(\mathcal H_9\) is injective. In particular,
\[
 \mathcal R_9(n+9)=\mathcal R_9(n),
 \qquad
 \mathcal H_9(n+9)\neq\mathcal H_9(n).
\]
\end{proposition}

\begin{proof}
The first equality follows from
\cref{p12-83-c1-s4:cor:retorno-fase-incremento-memoria}. The third coordinate of
\(\mathcal H_9(n+9)\) exceeds that of \(\mathcal H_9(n)\) by one unit, hence
both points are distinct. Finally, \(n\) is recovered from the third coordinate
by multiplying by nine and adding one to the phase part;
therefore, \(\mathcal H_9\) is injective.
\end{proof}

This proposition is the minimal model of the principle that will run through the entire
work:
\[
 \boxed{\text{phase return}\neq\text{state return}.}
\]
The nonadic connection of the Topological Phase Kernel will preserve more fields than
the elementary helix ---sheet, orientation, carry, boundary, terminal, and
route memory---, but it will not alter this inaugural distinction.

\subsection{Correspondence between Discrete Cardinality and Projective Measure}

\begin{definition}[Counting]
The count of a finite set \(X\) is its cardinality \(|X|\). It requires
neither an orientation, a scale, nor an additional reader.
\end{definition}

\begin{definition}[Measured data]
A discrete measurement datum is a triple
\[
 (x,\mathcal L,\mathcal C),
\]
where \(x\) is a state, \(\mathcal L\) is a typed reader, and
\(\mathcal C\) is a publication chart compatible with the codomain of
\(\mathcal L\). Its visible output is
\[
 \mathcal C(\mathcal L(x)).
\]
\end{definition}

The same set may have a single cardinality and admit multiple measures,
because the local cochain, the orientation of integration, or the
publication chart may vary. Conversely, two distinct states may publish the
same coordinate. The work will therefore always preserve two levels:
\[
 \text{state with genealogy}
 \quad\longrightarrow\quad
 \text{published coordinate}.
\]

\begin{criterion}[Refutation criteria for the inaugural unit]
The construction is refuted if any of the following
facts occurs:
\begin{enumerate}
 \item ten independent phases are counted in the residue cell;
 \item the mark \(9\) is identified with the integer \(0\) without declaring the
 change of chart;
 \item the passage \(9\to1\) forces \(q\mapsto0\);
 \item two distinct integers have the same pair \((q_9,\rho_9)\);
 \item a physical unit is attributed to the cell before defining the
 dimensional transducer.
\end{enumerate}
\end{criterion}

\subsection{Inaugural domain of the nonadic cell and causal exclusion of APP, TRIT, TPK, and analytic constants}

The unit constructed so far contains only:
\[
 \mathsf P_{10},\quad
 \mathsf D_9,\quad
 \lambda_9,\quad
 (q_9,\rho_9),\quad
 \mathcal R_9,\quad
 \mathcal H_9.
\]
It does not yet contain APP, TRIT, the Topological Phase Kernel, or an analytic
constant. The following unit will take two copies of the nonadic cell,
construct its gnomonic completion, and prove how the bidimensional support appears
over which the multiplicative matrix of APP and its internal additive law act.


% Unidad U003B. Esqueleto ciclotómico de APP y completación binomial.
% Redacción autónoma a partir de los objetos, no de la arquitectura del PDF abreviado.

\section{Cyclotomic skeleton of APP and binomial completion of the unit}
\label{p12-83-c1-s5:chap:app-esqueleto-completacion}

The defect between the APP matrix and its internal accumulator is not an isolated
number. Its law in
Cartesian length, its decomposition into orbits of order three, and the
completion of the nonadic cell produce three different objects that must
remain typed:
\[
 \text{scalar defect}\quad D_d,
 \qquad
 \text{cyclotomic module}\quad W_{\rm APP},
 \qquad
 \text{completion}\quad\overline{\mathcal C}_9.
\]
The first counts; the second preserves orientation and rank; the third adds
boundary strata. Equalities between their cardinals do not authorize identifying
their elements.

\subsection{Defect law at fixed modulus and variable length}

For each integer \(d\geq1\), let
\[
 B_d:=\{1,\ldots,9\}^d.
\]
The nonary digital root is represented by
\[
 \operatorname{dr}_9(n)=
 \begin{cases}
 n\bmod9,&n\not\equiv0\pmod9,\\
 9,&n\equiv0\pmod9.
 \end{cases}
\]
We define the additive accumulation and the multiplicative matrix evaluation on the same domain:
\[
 A_d(x_1,\ldots,x_d)
 =\operatorname{dr}_9\!\left(\sum_{j=1}^d x_j\right),
\]
\[
 M_d(x_1,\ldots,x_d)
 =\operatorname{dr}_9\!\left(\prod_{j=1}^d x_j\right).
\]
Their visible masses and their defect are
\[
 S_\Sigma(d):=\sum_{x\in B_d}A_d(x),
 \qquad
 S_\Pi(d):=\sum_{x\in B_d}M_d(x),
\]
\[
 D_d:=S_\Pi(d)-S_\Sigma(d).
\]
Here \(d\) counts Cartesian positions. The module \(9\) does not vary, and no
physical dimension has yet been introduced.

\begin{lemma}[Mass of the additive accumulator]
\label{p12-83-c1-s5:lem:masa-hoja-aditiva-longitud}
For every \(d\geq1\),
\[
 S_\Sigma(d)=45\,9^{d-1}.
\]
\end{lemma}

\begin{proof}
Once the first \(d-1\) coordinates and a published residue
\(r\in\{1,\ldots,9\}\) have been fixed, there exists a unique last coordinate that makes the
sum have residue \(r\). Each visible value thus has \(9^{d-1}\)
preimages. Summing the weights \(1+\cdots+9=45\) gives the formula.
\end{proof}

\begin{lemma}[Census of the multiplicative matrix]
\label{p12-83-c1-s5:lem:censo-hoja-multiplicativa-longitud}
For each of the six unit residues there are \(6^{d-1}\) preimages for each
fixed product value. The residues \(3\) and \(6\) have
\(d6^{d-1}\) preimages each. The zero residue, published as
\(9\), has
\[
 9^d-6^d-2d6^{d-1}
\]
preimages. Consequently,
\[
 S_\Pi(d)=9^{d+1}-9(d+3)6^{d-1}.
\]
\end{lemma}

\begin{proof}
The unit group of \(\mathbb Z/9\mathbb Z\) has six elements. Upon
fixing \(d-1\) units, the last is determined by the prescribed
unit product, giving \(6^{d-1}\) solutions for each of the
six classes.

To obtain a product of ternary valuation exactly one, one chooses which
of the \(d\) coordinates belongs to one of the two classes \(3,6\), while
the others are units. Once the product \(3\) or \(6\) is fixed, the last unit
is determined; hence \(d6^{d-1}\) for each class. All remaining tuples
have product zero modulo \(9\). Subtracting the two preceding
counts from \(9^d\) gives their multiplicity. The sum weighted by the
visible representatives yields the final identity.
\end{proof}

\begin{theorem}[Internal dimensional law of APP]
\label{p12-83-c1-s5:thm:ley-dimensional-app-autonoma}
For every \(d\geq1\),
\[
 \boxed{D_d=4\,9^d-9(d+3)6^{d-1}.}
\]
In particular,
\[
 D_1=0,
 \qquad
 D_2=54,
\]
and
\[
 \sum_{d\geq1}D_dz^d
 =\frac{54z^2(1-3z)}{(1-9z)(1-6z)^2}.
\]
Moreover,
\[
 \frac{D_d}{9^d}
 =4-(d+3)\left(\frac23\right)^{d-1}
 \longrightarrow4.
\]
\end{theorem}

\begin{proof}
The formulas from the two lemmas are subtracted. For the generating function, one uses
\[
 \sum_{d\geq1}9^dz^d=\frac{9z}{1-9z},
 \qquad
 \sum_{d\geq1}(d+3)6^{d-1}z^d
 =\frac{z(4-18z)}{(1-6z)^2},
\]
and the common numerator is reduced. The normalized formula follows by dividing
by \(9^d\); the exponential term multiplied by \(d+3\) converges to
zero.
\end{proof}

The equality \(D_2=54\) reproduce the census \(459-405\), but the formula
before and the law of modulo variable
\(\Delta_{3^m}^{\rm pl}=m3^{2m-1}\) are enunciados distintos: the first
varia \(d\) manteniendo the modulo \(9\); the second varia the modulo.

\subsection{Primary cyclotomic operator}

Let
\[
 G_{A_2}=\begin{pmatrix}2&-1\\-1&2\end{pmatrix},
 \qquad
 C=\begin{pmatrix}0&-1\\1&-1\end{pmatrix}.
\]

\begin{proposition}[Coxeter of type \(A_2\)]
\label{p12-83-c1-s5:prop:coxeter-a2-autonomo}
They are verified.
\[
 C^3=I_2,
 \qquad
 C^{\mathsf T}G_{A_2}C=G_{A_2},
 \qquad
 \det(XI_2-C)=X^2+X+1.
\]
Therefore \(C\) preserves the form \(A_2\), has order exactly three, and
has no nonzero fixed vector.
\end{proposition}

\begin{proof}
The multiplication gives
\[
 C^2=\begin{pmatrix}-1&1\\-1&0\end{pmatrix},
 \qquad C^3=I_2.
\]
Direct substitution into \(C^{\mathsf T}G_{A_2}C\) yields
\(G_{A_2}\). Finally, \(1\) is not a root of \(X^2+X+1\), hence
\(\ker(C-I_2)=0\).
\end{proof}

\subsection{Decomposition of APP into \(12\) fibers of type \(A_2\)}

Consider the action
\[
 a:\mathbb Z/9\mathbb Z\longrightarrow\mathbb Z/9\mathbb Z,
 \qquad a(r)=4r.
\]
Since \(4^3\equiv1\pmod9\), its restriction to the units has two free
orbits:
\[
 \mathcal O_1=(1,4,7),
 \qquad
 \mathcal O_2=(2,8,5).
\]
For an orbit \(\mathcal O=(r,4r,7r)\), we define the augmentation module
\[
 A(\mathcal O):=
 \ker\!\left(\mathbb Z[\mathcal O]
 \xrightarrow{\,\sum\,}\mathbb Z\right).
\]
At the base
\[
 b_1=e_r-e_{4r},
 \qquad b_2=e_{4r}-e_{7r},
\]
the restricted Euclidean form has Gram matrix \(G_{A_2}\), and the cycle
\(r\mapsto4r\mapsto7r\mapsto r\) acts by \(C\).

For \(x=(x_1,x_2,x_3)\in(\mathbb Z/9\mathbb Z)^3\), the three unordered pairs are
\[
 \{1,2\},\quad\{2,3\},\quad\{3,1\}.
\]
Over each pair, two typed laws are evaluated:
\[
 \ell_{ij}^{\Sigma}(x)=x_i+x_j,
 \qquad
 \ell_{ij}^{\Pi}(x)=x_ix_j.
\]
Since
\[
 \ell_{ij}^{\Sigma}(ax)=a\ell_{ij}^{\Sigma}(x),
 \qquad
 \ell_{ij}^{\Pi}(ax)=a^{-1}\ell_{ij}^{\Pi}(x),
\]
the sum carries the orientation \(C\) and the product the inverse orientation
\(C^{-1}\).

\begin{definition}[Cyclotomic module of APP]
\label{p12-83-c1-s5:def:modulo-ciclotomico-app-autonomo}
It is defined
\[
 \boxed{
 W_{\rm APP}:=
 \bigoplus_{\{i,j\}}
 \bigoplus_{\mu\in\{\Sigma,\Pi\}}
 \bigoplus_{\mathcal O\in\{\mathcal O_1,\mathcal O_2\}}
 A(\mathcal O).}
\]
There are \(3\cdot2\cdot2=12\) summands, each of rank two; therefore
\[
 W_{\rm APP}\cong A_2^{12},
 \qquad \operatorname{rank}W_{\rm APP}=24.
\]
The par--law--orbit labels are part of the object.
\end{definition}

For \(u\in\mathbb Z/9\mathbb Z\), we write
\[
 \beta_{\mathcal O}(u)=
 \begin{cases}
 (1,0),&u=r,\\
 (0,1),&u=4r,\\
 (-1,-1),&u=7r,\\
 (0,0),&u\notin\mathcal O,
 \end{cases}
\]
in the basis \((b_1,b_2)\). The state map is
\[
 \Psi(x)_{ij,\mu,\mathcal O}
 :=\beta_{\mathcal O}(\ell_{ij}^{\mu}(x)).
\]

\begin{theorem}[Equivariance and rank of the state skeleton]
\label{p12-83-c1-s5:thm:esqueleto-estatal-app-autonomo}
Let \(\rho(a)\) be the action using \(C\) on the six additive
fibers and \(C^{-1}\) on the six multiplicative fibers. Then
\[
 \boxed{\Psi(ax)=\rho(a)\Psi(x)}
\]
for the \(729\) states of \((\mathbb Z/9\mathbb Z)^3\). The rational submodule
generated by the images has rank \(24\).
\end{theorem}

\begin{proof}
Under the additive law, multiplying all coordinates by \(4\) advances one
position in each orbit; under the multiplicative law, it multiplies the value by
\(4^2\equiv7\equiv4^{-1}\pmod9\) and therefore reverses the cycle. This proves
componentwise equivariance.

In order for the rank not to remain hidden in a computational statement, we fix the coordinate order.
\[
 (12,\Sigma,\mathcal O_1),(12,\Sigma,\mathcal O_2),
 (12,\Pi,\mathcal O_1),(12,\Pi,\mathcal O_2),
\]
followed by the same order for the pairs \(23\) and \(31\), and, within each
fiber, the two coordinates \((b_1,b_2)\). Consider the states
\[
\begin{gathered}
001,002,004,005,010,011,012,013,014,015,016,020,\\
022,023,040,050,101,102,104,105,110,120,140,150.
\end{gathered}
\]
The integer matrix \(M_\Psi\) whose rows are its twenty-four images has
entries in ({-1,0,1}). Bareiss's fraction-free elimination yields the following
successive leading principal minors
\[
1,1,1,-1,-1,-1,-1,1,-1,1,1,1,1,-2,-1,1,1,1,2,1,4,-6,12,9.
\]
In particular,
\[
 \det M_\Psi=9\neq0.
\]
These twenty-four images are linearly independent. Since the
codomain already has rank \(24\), they generate \(W_{\rm APP}\otimes\mathbb Q\).
\end{proof}

The proof exhibits a finite reproducible certificate; it does not deduce full rank
from the mere surjectivity of each component projection.

\subsection{Localization of the aggregate \texorpdfstring{\(54\)}{54}}

Let
\[
 \mathcal M:=\mathbb Z[\mathbb Z/9\mathbb Z]
\]
and, for \(\mu\in\{\Sigma,\Pi\}\),
\[
 \nu_\mu:=\sum_{x,y\in\mathbb Z/9\mathbb Z}
 e_{\ell^\mu(x,y)}.
\]
The census of the two tables gives
\[
 \delta_{\Pi\Sigma}:=\nu_\Pi-\nu_\Sigma
 =12e_0+3e_3+3e_6
 -3\sum_{u\in(\mathbb Z/9\mathbb Z)^\times}e_u.
\]
Its coefficients are constant on the orbits of \(a\); therefore
\[
 (I-a)\delta_{\Pi\Sigma}=0.
\]
We define the increase projection
\[
 T:\mathcal M\longrightarrow
 A(\mathcal O_1)\oplus A(\mathcal O_2),
\]
\[
 T(\nu)=
 \left(((I-a)\nu)|_{\mathcal O_1},
       ((I-a)\nu)|_{\mathcal O_2}\right).
\]
Then
\[
 T(\delta_{\Pi\Sigma})=0.
\]
However, the visible reader
\[
 w(e_0)=9,
 \qquad w(e_r)=r\quad(1\leq r\leq8)
\]
satisfies
\[
 w(\delta_{\Pi\Sigma})=54.
\]

\begin{corollary}[Isotopic Obstruction]
\label{p12-83-c1-s5:cor:obstruccion-isotipica-app-autonoma}
The scalar \(54\) belongs to the \(C_3\)-invariant sector and cannot
be recovered by means of a linear functional that factors exclusively through
the augmentation modules. In particular,
\[
 \sum_x\Psi(x)=0
\]
and, for every linear functional \(L\) on \(W_{\rm APP}\otimes\mathbb Q\),
\[
 L\!\left(\sum_x\Psi(x)\right)=0\neq54.
\]
\end{corollary}

This obstruction prevents presenting \(54\) as a retrospective linear sum
of the twelve fibers. The subsequent APP--Fock coupling preserves the
orientation and uses an additional object; it does not contradict the corollary.

\subsection{Cubic completion of the nonad cell}

Let \(E_9\) be the underlying set of \(\mathbb Z/9\mathbb Z\), and let
\(\star\notin E_9\) be a boundary state. We define
\[
 \mathcal C_9:=E_9^3,
 \qquad
 \overline{\mathcal C}_9:=(E_9\sqcup\{\star\})^3.
\]
The symbol \(\star\) is not the zero class: it indicates that a coordinate has
reached the closure stratum.

\begin{theorem}[Binomial stratification of the unit cube]
\label{p12-83-c1-s5:thm:emrh-binomial-autonoma}
The stratum with exactly \(r\) coordinates equal to \(\star\) contains
\[
 \binom3r9^{3-r}
\]
elements. Therefore,
\[
 \boxed{1000=729+243+27+1=729+271.}
\]
Upon removing only the apex \((\star,\star,\star)\), there remains
\[
 999=729+243+27=729+270.
\]
\end{theorem}

\begin{proof}
The \(r\) boundary coordinates are chosen in \(\binom3r\) ways; each
remaining coordinate admits nine values. The four strata are disjoint
and exhaust the Cartesian product. The sum is the expansion of
\((9+1)^3\).
\end{proof}

\begin{figure}[htbp]
\centering
\resizebox{0.96\linewidth}{!}{\input{figuras/base_autoridad/fig_cubo_emrh_autoridad}}
\caption{Cubic completion of the nonadic cell. The interior of
cardinality \(729\) and the boundary strata of cardinalities \(243\), \(27\), and
\(1\) are disjoint parts of a single object.}
\label{p12-83-c1-s5:fig:cubo-emrh-autonomo}
\end{figure}

This stratification receives in HMT the name of extreme and mean
holographic ratio. Mathematically it is the typed binomial completion of the nonadic cell;
the name does not replace its definition.

For every \(d\geq1\), let
\[
 \overline{\mathcal C}^{(d)}_9=(E_9\sqcup\{\star\})^d.
\]
Asignemos peso \(1/10\) to each symbol of a coordinate and tomemos the measure
product \(\mu_d\).

\begin{theorem}[Projective compatibility of the measure]
\label{p12-83-c1-s5:thm:medida-proyectiva-completacion}
The projection that forgets the last coordinate satisfies
\[
 (\pi_d)_*\mu_d=\mu_{d-1}.
\]
Moreover,
\[
 10^d=\sum_{r=0}^{d}\binom dr9^{d-r},
 \qquad
 \mu_d(\overline{\mathcal C}^{(d)}_9)=1.
\]
The number of states increases with \(d\), while the normalized total mass
is conserved.
\end{theorem}

\begin{proof}
For every cylinder \(A\times(E_9\sqcup\{\star\})\),
\[
 \mu_d(A\times(E_9\sqcup\{\star\}))
 =\mu_{d-1}(A)\sum_{u\in E_9\sqcup\{\star\}}\frac1{10}
 =\mu_{d-1}(A).
\]
The cilindros generate the algebra finite of subconjuntos, what that determines the
measure image. The identity cardinal again becomes the binomio of Newton.
\end{proof}

\subsection{Morphism from the central block of APP to its Euclidean realization}

The already constructed central block is
\[
 B_c=\begin{pmatrix}7&2\\2&7\end{pmatrix}.
\]
Row reading and diagonal reading give
\[
 72+27=77+22=99,
\]
\[
 72-27=45=\det B_c,
 \qquad
 77-22=55=54+1.
\]
The relation with the completion is not obtained by concatenating symbols, but by the Euclidean division.
\[
 729=10\cdot72+9,
 \qquad
 271=10\cdot27+1.
\]

\begin{theorem}[Uniqueness of the Decimal Lift]
\label{p12-83-c1-s5:thm:levantamiento-decimal-emrh-autonomo}
For \(L_d(n)=10n+d\), with \(d\in\{0,\ldots,9\}\), there exists a single pair
\((d,d')\) such that
\[
 L_d(72)=729,
 \qquad
 L_d(72)+L_{d'}(27)=1000,
 \qquad
 d+d'=10.
\]
It is the pair
\[
 (d,d')=(9,1).
\]
\end{theorem}

\begin{proof}
The first equality is \(720+d=729\), hence \(d=9\). The second becomes
\(729+270+d'=1000\), whence \(d'=1\). The third is verified, and the
uniqueness of quotient and remainder rules out any other solution.
\end{proof}

Thus the causal diagram is obtained.
\[
 B_c\longrightarrow(72,27)
 \longleftarrow((72,9),(27,1))
 \longleftarrow(729,271).
\]
The block and its reading belong to APP; the dividends and remainders belong
to the completion. The diagram shares quotients; it does not invert
causality.

\subsection{Incidence of the cube and separation with respect to sextets and octets}

Let \(M_d\) be the incidence matrix of the \(2d\) codimension-one faces
of a hypercube on its \(2^d\) vertices. Each face contains
\(2^{d-1}\) vertices; two opposite faces are disjoint, and two
transversal faces share \(2^{d-2}\) vertices.

\begin{theorem}[Hypercube incidence spectrum]
\label{p12-83-c1-s5:thm:espectro-incidencia-cubo-autonomo}
\[
 \operatorname{Spec}(M_dM_d^{\mathsf T})
 =\left\{
 d2^{d-1}\,^{(1)},
 2^{d-1}\,^{(d)},
 0^{(d-1)}
 \right\}.
\]
Para \(d=3\),
\[
 \operatorname{Spec}(M_3M_3^{\mathsf T})
 =\{12,4,4,4,0,0\}.
\]
\end{theorem}

\begin{proof}
The row space decomposes into: the constant line; the \(d\)
differences of each pair of opposite faces; and the (d-1) combinations of
pair sums whose total sum is zero. The intersections described above
cause the Gram matrix to act by \(d2^{d-1}\), \(2^{d-1}\), and \(0\),
respectively.
\end{proof}

The three-dimensional cube has six faces, eight vertices, and twelve edges. Those
cardinals anticipate an incidence pattern (6/8), but a face is not
a hexad and the set of vertices is not an octad. The exceptional
identification requires a subsequent transporter that preserves supports and
gauge.

\subsection{Resonance of the incidence operator in rank \(6\)}

Define the strata that preserve exactly \(k\) nonadic
coordinates in rank \(d\):
\[
 \mathcal I_k(d):=\binom dk9^k.
\]
The defects of APP are
\[
 D_{\rm vis}=54,
 \qquad
 D_{\rm tot}=1215,
 \qquad
 D_{\rm coc}=129,
 \qquad
 1215=54+9\cdot129.
\]

\begin{theorem}[Uniqueness of the Binomial Resonance]
\label{p12-83-c1-s5:thm:resonancia-binomial-autonoma}
The unique positive integer \(d\) that simultaneously satisfies
\[
 \mathcal I_1(d)=54,
 \qquad
 \mathcal I_2(d)=1215
\]
is \(d=6\). In that range,
\[
 \frac{\mathcal I_2(6)-\mathcal I_1(6)}9=129.
\]
\end{theorem}

\begin{proof}
The first equality is \(9d=54\), which forces \(d=6\). Then
\[
 \mathcal I_2(6)=81\binom62=81\cdot15=1215,
\]
and
\[
 \frac{1215-54}{9}=129.
\]
There remains no other possible rank because the first equality already determines \(d\).
\end{proof}

\subsection{Obstructions of the cyclotomic skeleton of APP: fibers \(A_2^{12}\), completion \(729+271\), rank resonance \(6\), and cubic incidence}

\begin{criterion}[Typed falsification of the unit]
The construction is refuted if any of the following facts occurs:
\begin{enumerate}
 \item the law of \(D_d\) not reproduce \(D_1=0\) and \(D_2=54\);
 \item variation of length is fused with variation of modulus;
 \item \(C\) ceases to preserve \(G_{A_2}\) or acquires a fixed vector;
 \item the labels pair--sheet--orbit not produce twelve fibers;
 \item the smallest certificate of \(\Psi\) has zero determinant;
 \item a functional that factors through the augmentation module recovers \(54\);
 \item the symbol of boundary \(\star\) it identifica with the zero residue;
 \item the four strata do not sum to \(1000\);
 \item a cubic face is identified with a hexad without an incidence
       transporter.
\end{enumerate}
\end{criterion}

The unit has constructed, without recognizing them retrospectively, the module
\(A_2^{12}\), the completion \(729+271\), the projective measure, and the
rank-six resonance. These structures will feed, in distinct ways,
the generation of constants and the exceptional branches.


% Unidad U006F. Inventario operatorio completo del Topological Phase Kernel.

\long\def\HMTDeferredTPKDevelopment{%
\section{Operator algebra of the Topological Phase Kernel and causal update dynamics of the enriched state}
\label{p12-83-c1-s6:chap:inventario-operatorio-completo-tpk}

The Topological Phase Kernel is neither a reduction modulo three nor a clock with one hundred and eight positions. It is a category of enriched states equipped with operators of selection, transport, reading, memory, periodicity, prolongation, and publication. To prevent an abbreviation from replacing the object, each operator is recorded by its type and by the information it preserves.

\subsection{Typed schema of the 34 canonical TPK constructions: domain, codomain, action, invariants, and memory}

\begin{definition}[Operator specification]
\label{p12-83-c1-s6:def:u006f-ficha-operatoria}
A canonical realization of the TPK is a sextuple
\begin{equation}
 \mathcal O=
 (D_{\mathcal O},C_{\mathcal O},f_{\mathcal O},
 I_{\mathcal O},L_{\mathcal O},M_{\mathcal O}),
 \label{p12-83-c1-s6:eq:u006f-ficha}
\end{equation}
where:
\begin{enumerate}
 \item \(f_{\mathcal O}:D_{\mathcal O}\to C_{\mathcal O}\) is its rule;
 \item \(I_{\mathcal O}\) is the family of conserved relations;
 \item \(L_{\mathcal O}\) declares the information published or lost;
 \item \(M_{\mathcal O}\) enumerates the memory that remains in the enriched
 state.
\end{enumerate}
\end{definition}

Two realizations represent the same operator only when there exists a
typed equivalence that intertwines rules and invariants. Sharing a period,
cardinality, or visible output is not sufficient.

\subsection{Classification of the enriched state into \(8\) invariant structural classes}

The inventory is partitioned into:
\begin{enumerate}
 \item[\textbf{C1.}] supports and states;
 \item[\textbf{C2.}] internal selection;
 \item[\textbf{C3.}] transporte;
 \item[\textbf{C4.}] readers and quotients;
 \item[\textbf{C5.}] memory and boundary;
 \item[\textbf{C6.}] periodicity and return;
 \item[\textbf{C7.}] prolongation;
 \item[\textbf{C8.}] salidas.
\end{enumerate}
The class is determined by domain, codomain, and function, not by the name of a number associated with it.

\Needspace{.92\textheight}
\subsection{Functional decomposition into \(14\) groupings compatible with the
APP–TRIT–TPK operators}

\begingroup
\small
\renewcommand{\arraystretch}{1.12}
\begin{longtable}{@{}>{\raggedright\arraybackslash}p{.20\textwidth}>{\raggedright\arraybackslash}p{.29\textwidth}>{\raggedright\arraybackslash}p{.43\textwidth}@{}}
\toprule
Grouping&Operators or spaces&Function and invariants\\
\midrule
\endfirsthead
\toprule
Grouping&Operators or spaces&Function and invariants\\
\midrule
\endhead
State and route memory&
\(\mathcal X_{\rm TPK}^{\rm enr}\), \(\mathcal F_q\)&
Preserve position, sheet, orientation, residue, quotient, carry, signature,
boundary, cylinder, prefix, provenance and memory.\\

Internal selection&
\(M_{\rm ph}:\{1,\ldots,9\}\to\{-1,0,+1\}\)&
Activate the additive, multiplicative evaluation or the threshold step; does not move
the state.\\

Transporte espacial&
\(T_N,T_E,T_S,T_O,F_{\rm loc},F_{\rm obs}\)&
Perform cardinal translations, sheet feedback, and timeline of two cursors.\\

Lectores modulares&
\(\rho_9,q_9,\rho_3^{\rm or},c_3,q_{1000},r_{1000}\)&
Publish phase, orientation and decimal block preserving the lift quotients.\\

Block and memory&
\(\mathcal K_{27}=F_{\rm obs}^{27}\), \(\mathcal N_{27}\)&
The former composes twenty-seven transports; the latter filters memory
events. An equal index does not imply an equal operator.\\

Twelve-phase state&
\(\mathcal R_{12},C_{12},\Theta_{108},\Gamma_{108}\)&
Distinguish enriched record, sector, cyclic law, and ordered support.\\

Periodicity and depths&
\(C_{12}\hookrightarrow C_{108}\twoheadrightarrow C_9\),
\(w_{108},w_{1080}\)&
Preserves the non-splittable extent, words of different lengths and
the vertical memory.\\

Channels \(90/120\)&
\(\mathcal F_6,\mathcal F_8,\mathcal I_{6\to8},
S_{90},S_{120}\)&
Separate finite incidences, analytic tails, and their normalizations.\\

Bidirectionality&
\(P_\pm,\operatorname{rev},\operatorname{Ext}_\pm,
N_\pm,\beta,C_M\)&
Perform exchange, route inversion, future and past prolongation, valence, balance, and conjugation.\\

Emission and coinduction&
\(L_0,L_1,R_{36},G_9,\varprojlim\operatorname{Cyl}_m\)&
Lift the seed, open the boundary, and organize finite and compatible extensions.\\

Numeric Readers&
\(\mathscr C_\pi,\mathscr P_e,F_{\rm av}\)&
They construct rational forks, factorial propagation, and Fibonacci self-scaling.\\

Exceptional Projection&
\(\Gamma_{A_W}\), weight, projective support&
Transport the same emission to ternary incidence without selecting Archimedean numbers.\\

Atlas and realization&
route signature, atlas \(81/56/13/324\), \(C_M\)&
Classify stories and characters before applying a dimensional chart.\\

Finite validation&
censuses, structural necessity tests, integrity identities and
independence controls&
Checks coverage at the executed depth and separates construction,
independent verification, and comparison data.\\
\bottomrule
\end{longtable}
\endgroup

The eight classes describe mathematical nature; the fourteen groupings describe
execution function. They are two gradings of the same inventory.

\Needspace{.92\textheight}
\subsection{Domain and typing of the \(34\) entries of the update operator}

\begin{equation}
 \mathfrak O_{\rm TPK}:=(\mathfrak o_1,\ldots,\mathfrak o_{34}).
 \label{p12-83-c1-s6:eq:u006f-registro-operadores}
\end{equation}

\begingroup
\scriptsize
\renewcommand{\arraystretch}{1.12}
\begin{longtable}{@{}r p{.19\textwidth}p{.28\textwidth}p{.42\textwidth}@{}}
\toprule
\#&Symbol&Domain and codomain&Rule or function\\
\midrule
\endfirsthead
\toprule
\#&Symbol&Domain and codomain&Rule or function\\
\midrule
\endhead
1&\(\mathcal A_9\)&\((\mathbb Z/9)^2\), matrix and internal law&
Additive-multiplicative arithmetic cell.\\
2&\((\rho_9,q_9)\)&\(\mathbb Z_{>0}\to
\{1,\ldots,9\}\times\mathbb Z_{\geq0}\)&
Nonadic division with positive representative and quotient.\\
3&\(\rho_3^{\rm or},c_3\)&\(\mathbb Z\to
\{-1,0,+1\}\times\mathbb Z\)&
Lift ternary oriented \(n=\rho_3^{\rm or}(n)+3c_3(n)\).\\
4&\(\iota_{1000}:=(q_{1000},r_{1000})\)&\(\mathbb Z\to
\mathbb Z\times\{0,\ldots,999\}\)&
Euclidean division \(N=1000q_{1000}+r_{1000}\).\\
5&\(T_d\)&\(X_9\to X_9\), \(d\in\{N,E,S,O\}\)&
Cardinal oriented translation.\\
6&\(F_{\rm loc}\)&\(\mathcal Q_{\rm loc}\times
\mathscr D_{\rm card}\to\mathcal Q_{\rm loc}\)&
Local position feedback and sheet.\\
7&\(T_{\min}\)&\(\Omega_{\rm seed}\to\mathcal U_6\)&
Minimum emitter of the two cursors, including the coordinate signature.\\
8&\(G\)&\(\operatorname{im}s_\Sigma\times
\operatorname{im}s_\Pi\to\{0,\ldots,999\}^6\)&
Coordinate-by-coordinate decimal coupling.\\
9&\(\Psi_6\)&\(\mathcal U_6\to\mathbb F_3^6\)&
Tertiary reduction of the catalog.\\
10&\(\mathcal R_3\)&\(\mathcal Q_{\rm loc}\times
\mathscr D_{\rm card}^3\to\{-1,0,+1\}\)&
Local three-step observable factor.\\
11&\(\mathcal K_{27}\)&\(\mathcal Q_{\rm obs}\to\mathcal Q_{\rm obs}\)&
Composition \(F_{\rm obs}^{27}\).\\
12&\(\mathcal X_{\rm TPK}^{\rm enr}\)&enriched states category&
Preserve all the typed coordinates of the nucleus.\\
13&\(\mathcal H\)&relation on enriched states&
Hensel Transition that retains residue and quotient.\\
14&\(\mathcal C_{3,1000}\)&crossed states
\((N,L,D,K,A,B)\)&
Exact update of ternary-decimal cylinders.\\
15&\(\Sigma_{\rm B}\)&boundary memories \(\to\) signature joint&
Publish the block orientation preserving its genealogical record.\\
16&\({\rm EF}\text{--}G_9\)&extension relation&
Filter compatible extensions and organize the nonadic genealogy.\\
17&\(\Gamma_9\)&enriched histories \(\longrightarrow\) histories&
Holonomy of a phase twist with boundary memory.\\
18&\(X_\chi\)&\(\varprojlim_m\operatorname{Surv}^{(\chi)}_m\)&
Global section when non-emptiness, compatibility, and uniqueness are demonstrated in full depth.\\
19&\(\mathcal P_{\rm APP}^{(2)}\)&bivariate path category&
Transports simultaneously the multiplicative evaluation and the additive accumulation.\\
20&\(F_{\rm obs}\)&\(\mathcal Q_{\rm obs}\to\mathcal Q_{\rm obs}\)&
Deterministic observable chronology with two cursors.\\
21&\(\mathcal A_{42}\)&alphabet of forty-two triads&
Extended decimal support of finite emission.\\
22&\(\operatorname{Ext}_r\)&
\(\operatorname{Ext}_r\subseteq\mathcal F_q\times\mathcal F_{q'}\)&
Typed relation of fiber prolongation.\\
23&\(\mathcal K_{\rm ph}\)&\(\mathcal U_6\times C_9\) about itself&
Transfer of phase-compatible explicit continuation.\\
24&\((P,H,Q)\)&\(\mathcal X_\times\to\mathcal X_\times\)&
Advance, half-turn, and quarter-turn of seeds.\\
25&\((c_{\rm mode},c_{\rm or})\)&caminos \(\to\mathbb Z^2\)&
Mode and orientation change cocycles.\\
26&\(E_{20}\)&finite states \(\to\) twenty blocks&
Finite triadic prolongation with sourcebook.\\
27&\(\mathcal R_{12}\)&history enriched \(\to\) record of twelve
sectors&
Timed system with route, signature, carry, and genealogical record.\\
28&\(R_{\rm sgn}\)&terminal state \(\to U_{12}^{\rm sgn}
\subseteq\mathbb Z^{12}\)&
Reader of the signed harmonic observable.\\
29&\(\mathcal H_{12}\)&\(\mathbb Z^{12}\to\mathbb Z^{12}\)&
Transformation by three orbits of step three; inverse
\(\mathcal H_{12}/4\) over the pertinent image.\\
30&\(\mathcal E_{\rm dod}=(D_3,D_4,Q)\)&
\(\mathbb Z^{12}\to\mathbb Z^{25}\)&
Transversal extractor set of arities three and four.\\
31&\(\mathcal C_\alpha\)&blocks and carries \(\to\) blocks&
Dodecaphasic closure recurrence in base thousand.\\
32&\((L_{\rm tick},\chi_{\rm mass})\)&enriched history \(\to\)
temporal character&
Posterior dimensional interface; does not participate in the numeric constructor of
this work.\\
33&\(\mathcal W_\pm\)&bidirectional lattice of index twelve&
Reading in both directions preserving orientation and origin.\\
34&\(\mathcal A_{\rm species}\)&extension families \(\to\) atlas&
Posterior realization interface, not a primitive fourth.\\
\bottomrule
\end{longtable}
\endgroup

\subsection{Causal composition of the enriched-state update}

\begin{definition}[Complete canonical state]
\label{p12-83-c1-s6:def:u006f-esquema-estado}
Let
\[
 \mathcal E:=
 \mathbb Z_{\rm quo}\times\mathbb Z_{\rm car}\times
 \mathsf{Mem}\times\mathsf{Term}\times\mathsf{Cyl}\times
 \mathsf{Fr}\times\mathsf{Pref}\times\Lambda .
\]
The canonical enriched state at instant \(t\) is
\begin{equation}
 x_t=(p_t,m_t,\sigma_t,b_t,r_t;
 q_t,c_t,\eta_t,s_t,C_t,\partial_t,N_t,\lambda_t)
 \in
 X_9\times\{\Sigma,\Pi\}\times\mathsf{TRIT}\times
 C_{12}\times C_9\times\mathcal E.
 \label{p12-83-c1-s6:eq:u006f-esquema-estado}
\end{equation}
The state abbreviations used in other units are projections of
this tuple. In particular, the quotient \(q_t\), carry
\(c_t\), terminal \(s_t\), cylinder \(C_t\), boundary
\(\partial_t\), prefix \(N_t\), and ledger \(\lambda_t\) are not omitted.
\end{definition}

Let \(\rho_t\) be the observable arrow determined by \(x_t\). The transport
laws of U005 and the transition \(\mathcal C_{3,1000}\) determine a memory
endomorphism
\begin{align}
 \mathbf E_t:\mathcal E&\longrightarrow\mathcal E,\notag\\
 (q,c,\eta,s,C,\partial,N,\lambda)&\longmapsto
 \left(
 \begin{aligned}
 &\mathsf H(\rho_t)q,
  \mathsf C(\rho_t)c+\kappa(\rho_t),
  \mathsf M(\rho_t)\eta,\\
 &\mathsf T(\rho_t)s,
  \mathsf{Cyl}(\rho_t)C,
  \mathsf{Fr}(\rho_t)\partial,\\
 &\mathsf{Pref}(\rho_t)N,
  \Lambda(\rho_t)\lambda
 \end{aligned}
 \right).
 \label{p12-83-c1-s6:eq:u006f-actualizacion-memoria}
\end{align}
Each component published in \eqref{p12-83-c1-s6:eq:u006f-actualizacion-memoria} has the
domain indicated by its fiber; \(\kappa(\rho_t)\) is the carry cochain.

We define three endomorphisms of the canonical product. If
\(\sigma'=M_{\rm ph}(g(t))\) and \(m'=\mu(\sigma',m)\), then
\begin{align}
 \mathsf{Sel}_t(p,m,\sigma,b,r;e)
  &:=(p,m',\sigma',b,r;e),\label{p12-83-c1-s6:eq:u006f-selector-tipado}\\
 \mathsf{Tra}_t(p,m,\sigma,b,r;e)
  &:=(T_{d(t)}p,m,\sigma,b,r;e),\label{p12-83-c1-s6:eq:u006f-transporte-tipado}\\
 \mathsf{Upd}_t(p,m,\sigma,b,r;e)
  &:=(p,m,\sigma,b',r';\mathbf E_t(e)),
 \label{p12-83-c1-s6:eq:u006f-memoria-tipada}
\end{align}
where \((b',r')\) is the destination of \(\rho_t:(b,r)\to(b',r')\). The causal update is a composition of maps, not the composition of a map with a tritic value:
\begin{equation}
 \boxed{\mathcal U_t:=
 \mathsf{Upd}_t\circ\mathsf{Tra}_t\circ\mathsf{Sel}_t.}
 \label{p12-83-c1-s6:eq:u006f-tuberia}
\end{equation}

\begin{theorem}[Closure of causal update]
\label{p12-83-c1-s6:thm:u006f-cierre-tuberia}
The map \(\mathcal U_t\) is an endomorphism of the complete canonical
state. Readers, prolongations, and projections are applied after
\(\mathcal U_t\) only when their domain has been produced.
\end{theorem}

\begin{proof}
\(\mathsf{Sel}_t\) changes only the mode and the TRIT;
\(\mathsf{Tra}_t\) changes only the APP position; and \(\mathsf{Upd}_t\) applies the
typed product \eqref{p12-83-c1-s6:eq:u006f-actualizacion-memoria} and the phase arrow. All three codomains are the
same product of \cref{p12-83-c1-s6:def:u006f-esquema-estado}; their composition is therefore defined and
preserves all state coordinates.
\end{proof}

\subsection{Typed separation between transport, memory, calendar, extraction and return operators}

The inventory imposes
\begin{equation}
 \mathcal S_4\neq D_4\neq\operatorname{sd}_4,
 \qquad
 \mathcal R_{12}\neq C_{12}\neq\Theta_{108}\neq\Gamma_{108},
 \qquad
 F_{\rm obs}\neq\mathcal K_{\rm ph},
 \label{p12-83-c1-s6:eq:u006f-no-identificacion-uno}
\end{equation}
and
\begin{equation}
 U_{12}^{\rm cat}\neq U_{12}^{\rm sgn},
 \qquad
 \mathcal E_{\rm dod}\neq D_{\rm raw}^{90,120}
 \neq\mathcal K_{6\to8}.
 \label{p12-83-c1-s6:eq:u006f-no-identificacion-dos}
\end{equation}
They are type incompatibilities, not mere value inequalities.

\begin{criterion}[Audit of operatorial completeness]
\label{p12-83-c1-s6:crit:u006f-completitud}
An execution is operatorially complete only if:
\begin{enumerate}
 \item declares the inputs of \(\mathfrak O_{\rm TPK}\) used;
 \item is recorded the domain and the codomain of each composition;
 \item preserves the causal order of \eqref{p12-83-c1-s6:eq:u006f-tuberia};
 \item distinguishes dynamics, transfer, reading, prolongation, and output;
 \item permits reconstruction from \(\lambda_t\) of every block change and
 every carry;
 \item does not promote a finite check to an infinite section without
 proving non-emptiness, compatibility, and uniqueness at every depth.
\end{enumerate}
Any symbol used without a type, or any fusion prohibited by
\eqref{p12-83-c1-s6:eq:u006f-no-identificacion-uno} and
\eqref{p12-83-c1-s6:eq:u006f-no-identificacion-dos}, constitutes a verifiable failure.
\end{criterion}


% Unidad U006E. Separación tipada de los canales 90/120 y de los símbolos
% epsilon. La realización de incidencia se construye después de W12 y W24.

\section{Operator decomposition of incidence channels \texorpdfstring{\(90/120\)}{90/120}: temporal selector, nilpotent generator and angular component}
\label{p12-83-c1-s7:chap:canales-90-120-epsilon}

The integers \(90\) and \(120\) appear in several strata of the construction. Their
numerical coincidence does not establish an identity of objects. Likewise,
the letter epsilon denotes three constructions with incompatible domains. The
purpose of this unit is to fix their types before using them and to indicate
precisely which realizations require structures that have not yet been
constructed.

\subsection{Temporal, nilpotent, and angular operators: domains and codomains}

\subsubsection{Action of the tritic temporal selector over \(\{-1,0,+1\}\): opening, threshold, and return with memory}

The internal calendar selector is
\begin{equation}
 \varepsilon_9:\{1,\ldots,9\}
 \longrightarrow\{-1,0,+1\},
 \qquad
 \varepsilon_9(r)=
 \begin{cases}
  +1,&r\in\{1,4,7\},\\
  -1,&r\in\{2,5,8\},\\
  0,&r\in\{3,6,9\}.
 \end{cases}
 \label{p12-83-c1-s7:eq:u006e-epsilon-temporal}
\end{equation}
Its codomain is the oriented TRIT. It determines which sheet acts and does not itself shift any position.

\subsubsection{Nilpotent parabolic generator \(J_0^2=0\) and evolution at the temporal threshold}

The class
\begin{equation}
 \varepsilon_{\rm nil}
 \in
 \mathcal Q_{\rm n}:=
 \mathbb F_3[\varepsilon_{\rm nil}]
 /(\varepsilon_{\rm nil}^{\,2})
 \label{p12-83-c1-s7:eq:u006e-epsilon-nilpotente}
\end{equation}
satisfies
\begin{equation}
 \varepsilon_{\rm nil}\neq0,
 \qquad
 \varepsilon_{\rm nil}^{\,2}=0.
 \label{p12-83-c1-s7:eq:u006e-epsilon-nilpotente-relacion}
\end{equation}
It is an element of a degenerate quadratic algebra; it is not a temporal selector.

\subsubsection{Angular residue induced after temporal selection and orientation preservation}

The symbol
\begin{equation}
 \varepsilon_{\rm res}^{\circ}\in\mathbb R
 \label{p12-83-c1-s7:eq:u006e-epsilon-residual-tipo}
\end{equation}
designates, when its geometric parameters are introduced, the residue of an
angular identity. At this stage only its type remains reserved. It is not
obtained from \(\varepsilon_9\) or from
\(\varepsilon_{\rm nil}\), and it does not intervene in the construction of
\(\pi,e,\varphi\) or \(\alpha\).

\begin{proposition}[No identification of the epsilon symbols]
\label{p12-83-c1-s7:prop:u006e-no-identificacion-epsilon}
There is no well-typed equality between
\[
 \varepsilon_9,\qquad
 \varepsilon_{\rm nil},\qquad
 \varepsilon_{\rm res}^{\circ}.
\]
\end{proposition}

\begin{proof}
The first object is a function between finite sets; the second, a
nilpotent element of an algebra over \(\mathbb F_3\); the third, a
real scalar with angular unit. Their domains, codomains, and laws of
composition are distinct.
\end{proof}

\subsection{Distribution of the indices \(90/120\) among the \(6\) objects of the incidence system}

\begin{definition}[Dodecaphasic Extractor]
\label{p12-83-c1-s7:def:u006e-extractor-dodecafasico}
The integral extractor
\begin{equation}
 \mathcal E_{\rm dod}:=(D_3,D_4,Q):
 \mathbb Z^{12}\longrightarrow\mathbb Z^{25}
 \label{p12-83-c1-s7:eq:u006e-extractor-dodecafasico}
\end{equation}
reads differences of arities three and four and a sum coordinate. The
subscripts are not angular lengths.
\end{definition}

\begin{definition}[Raw difference of sequences]
\label{p12-83-c1-s7:def:u006e-diferencia-cruda}
For \(|q|<1\), define
\begin{equation}
 S_m^{\rm raw}(q):=\frac{q^m}{1-q^{3m}},
 \qquad
 D_{\rm raw}^{90,120}(q)
 :=S_{90}^{\rm raw}(q)-S_{120}^{\rm raw}(q).
 \label{p12-83-c1-s7:eq:u006e-series-crudas}
\end{equation}
This object is an analytic function of \(q\); it is not a finite census.
\end{definition}

\begin{definition}[Marked families of incidence]
\label{p12-83-c1-s7:def:u006e-familias-incidencia}
Once a hexad \(H\), its image
\(\ell(H)\) in a binary dodecad \(D\), and the lifted octad
\(O_H\) have been constructed, one will define
\begin{align}
 \mathcal F_6(H)
 &:=\ell(H)\times\binom{\ell(H)}2,
 \label{p12-83-c1-s7:eq:u006e-familia-seis}\\
 \mathcal F_8(H)
 &:=O_H\times\binom{\ell(H)}2.
 \label{p12-83-c1-s7:eq:u006e-familia-ocho}
\end{align}
By construction,
\begin{equation}
 |\mathcal F_6(H)|
 =6\binom62=90,
 \qquad
 |\mathcal F_8(H)|
 =8\binom62=120.
 \label{p12-83-c1-s7:eq:u006e-cardinales-incidencia}
\end{equation}
The inclusion \(\ell(H)\hookrightarrow O_H\) inducira
\begin{equation}
 \mathcal I_{6\to8}:
 \mathcal F_6(H)\hookrightarrow\mathcal F_8(H).
 \label{p12-83-c1-s7:eq:u006e-inclusion-incidencia}
\end{equation}
\end{definition}

The \cref{p12-83-c1-s7:eq:u006e-familia-seis,p12-83-c1-s7:eq:u006e-familia-ocho} fix the
type from now on. Their effective existence will be proved after constructing
\(C_W\), \(W_{12}\), the binary datum \(G_{24},D,\ell\), and \(W_{24}\).
That branch is not anticipated here.

\begin{definition}[Normalized Queue Combination]
\label{p12-83-c1-s7:def:u006e-combinacion-colas}
For parameters \(q_-,q_+\in(0,1)\), let
\begin{equation}
 S_m(q):=\frac{q^m}{1-q^{3m}},
 \qquad
 \Delta S_m:=S_m(q_-)-S_m(q_+).
 \label{p12-83-c1-s7:eq:u006e-colas-normalizadas}
\end{equation}
The functional
\begin{equation}
 \mathcal K_{6\to8}
 :=12\,\Delta S_{90}-\Delta S_{120}
 \label{p12-83-c1-s7:eq:u006e-funcional-seis-ocho}
\end{equation}
is an analytic combination of tails. Its coefficient twelve belongs to its
normalization and does not turn it into the record
\(\mathcal R_{12}\).
\end{definition}

\begin{definition}[Semigroup of Degrees and Moments]
\label{p12-83-c1-s7:def:u006e-semigrupo-momentos}
The additive semigroup generated by both degrees is
\begin{equation}
 \mathfrak S_{90,120}
 :=\langle90,120\rangle
 =\{90a+120b:a,b\in\mathbb N_0\}.
 \label{p12-83-c1-s7:eq:u006e-semigrupo}
\end{equation}
For \(n\in\mathfrak S_{90,120}\), the oriented moment is
\begin{equation}
 s_n:=q_-^n-q_+^n.
 \label{p12-83-c1-s7:eq:u006e-momento}
\end{equation}
\end{definition}

The moments are differential monomials indexed by a semigroup; the
series of \eqref{p12-83-c1-s7:eq:u006e-colas-normalizadas} are infinite sums of those
moments over specific progressions.

\subsection{Normalized incidence defect}

When the inclusion \(\mathcal I_{6\to8}\) has been constructed, its
difference of cardinalities will be
\[
 |\mathcal F_6(H)|-|\mathcal F_8(H)|=90-120=-30.
\]
If normalization is declared
\[
 24\binom62=360,
\]
we obtain
\begin{equation}
 \frac{90-120}{24\binom62}
 =\frac{6-8}{24}
 =-\frac1{12}.
 \label{p12-83-c1-s7:eq:u006e-defecto-normalizado}
\end{equation}
The inclusion determines \(-30\); the rational \(-1/12\) depends moreover of
the declared normalization. No it what must present as if the cardinality
by if alone fixed the denominator.

\subsection{Typing invariants}

\begin{theorem}[Separation of the six objects]
\label{p12-83-c1-s7:thm:u006e-separacion-seis-objetos}
The objects
\begin{equation}
 \mathcal E_{\rm dod},\quad
 D_{\rm raw}^{90,120},\quad
 \mathcal I_{6\to8},\quad
 \mathcal K_{6\to8},\quad
 \mathfrak S_{90,120},\quad
 \varepsilon_{\rm res}^{\circ}
 \label{p12-83-c1-s7:eq:u006e-seis-objetos}
\end{equation}
are mutually distinct by type: they are, respectively, an integer reader,
an analytic function, a finite morphism, a combination of series, an
index semigroup, and a real angular scalar.
\end{theorem}

\begin{proof}
The conclusion is obtained by comparing their domains and codomains:
\[
 \mathbb Z^{12}\to\mathbb Z^{25},\quad
 \mathbb D\to\mathbb C,\quad
 \mathcal F_6(H)\to\mathcal F_8(H),\quad
 (0,1)^2\to\mathbb R,\quad
 \mathbb N_0^2\to\mathbb N_0,\quad
 \{\ast\}\to\mathbb R.
\]
There are no two identical signatures.
\end{proof}

\begin{criterion}[Refutation criteria]
\label{p12-83-c1-s7:crit:u006e-refutacion}
The classification is refuted if any of its domains or
codomains is altered, if the census of
\eqref{p12-83-c1-s7:eq:u006e-cardinales-incidencia} does not reproduce \(90\) and \(120\), if
\(-1/12\) is attributed to the inclusion without declaring the normalization, or if
the angular residual is presented as an output of the constants generator of
this monograph.
\end{criterion}


% Unidad U006A. Factor local, observacion estroboscopica y transicion exacta
% de cilindros.

\section{Local transport factor and ternary-decimal transition}
\label{p12-83-c1-s8:chap:tpk-factor-local-transicion-cilindros}

This unit separates three operations involved in the transport of the
Topological Phase Kernel that do not share a domain: the cardinal realization
of a trit, the stroboscopic sampling of a sequence, and the exact update
of a pair of ternary--decimal cylinders. The first is a finite
dynamics; the second is a temporal reader; the third preserves the integers
that make it possible to decide compatibility without floating-point arithmetic.

\subsection{Local state and sheet feedback}

Let
\[
 X_9=(\mathbb Z/9\mathbb Z)^2,
 \qquad
 \mathscr D_{\rm card}=\{N,E,S,O\},
\]
with
\[
 \delta(N)=(-1,0),\quad \delta(E)=(0,1),\quad
 \delta(S)=(1,0),\quad \delta(O)=(0,-1).
\]
The visible evaluations of APP are
\[
 L_\Sigma(x,y)=\rho_9((x+1)+(y+1)),
 \qquad
 L_\Pi(x,y)=\rho_9((x+1)(y+1)),
\]
where \(\rho_9(n)=1+((n-1)\bmod9)\). We further define
\[
 \rho_3^{\rm or}([0])=0,
 \qquad
 \rho_3^{\rm or}([1])=+1,
 \qquad
 \rho_3^{\rm or}([2])=-1.
\]

\begin{definition}[Local factor of APP--TRIT]
\label{p12-83-c1-s8:def:u006a-factor-local}
The local state space is
\[
 \mathcal Q_{\rm loc}=X_9\times\{\Sigma,\Pi\},
 \qquad |\mathcal Q_{\rm loc}|=162.
\]
For \(q=(p,m)\), \(d\in\mathscr D_{\rm card}\), we set
\[
 p_d=p+\delta(d),
 \qquad
 a_d(q)=\rho_3^{\rm or}([L_m(p_d)]_3),
\]
and
\[
 \mu(a,m)=
 \begin{cases}
  \Sigma,&a=+1,\\
  \Pi,&a=-1,\\
  m,&a=0.
 \end{cases}
\]
The local transition is
\[
 F_d^{\rm loc}(p,m)=(p_d,\mu(a_d(p,m),m)).
\]
\end{definition}

The position it is modified before evaluating the sheet. The null trit preserves the
last active mode; not eliminates the sheet nor reinitializes the history.

For a cardinal word \(u=d_1d_2d_3\), let
\[
 q_j=F_{d_j}^{\rm loc}(q_{j-1}),\qquad q_0=q,
\]
and let us denote by
\[
 R_3(q,u)\in\{-1,0,+1\}
\]
the trit produced in the third displacement.

\begin{definition}[Remark on step three]
\label{p12-83-c1-s8:def:u006a-obs3}
On an oriented sequence \(a=(a_t)_{t\geq1}\) we define
\[
 \operatorname{Obs}_3(a)=(a_3,a_6,a_9,\ldots).
\]
The map \(R_3\) takes a state and a cardinal word of length three;
\(\operatorname{Obs}_3\) takes a sequence and publishes a subsequence. They are
operators of different types.
\end{definition}

\begin{theorem}[Exact surjectivity of the sampled factor]
\label{p12-83-c1-s8:thm:u006a-sobreyectividad-r3}
For every \(q\in\mathcal Q_{\rm loc}\) and every
\(b\in\{-1,0,+1\}\), the fiber
\[
 \mathcal R_3(q,b)
 =\{u\in\mathscr D_{\rm card}^3:R_3(q,u)=b\}
\]
is non-empty and satisfies
\[
 \boxed{8\leq |\mathcal R_3(q,b)|\leq40.}
\]
In particular, the \(162\cdot3=486\) state--output fibers are non-empty.
\end{theorem}

\begin{proof}
From each of the \(162\) states, the \(4^3=64\) cardinal words are
enumerated, and the recurrence of
\cref{p12-83-c1-s8:def:u006a-factor-local} is applied literally. The exhaustive census of the \(486\) pairs
\((q,b)\) has a minimum of eight, a maximum of forty, and no zero value. The
statement is therefore a finite identity on the defined system, not
a mixing hypothesis.
\end{proof}

\begin{corollary}[Realization of an arbitrary oriented sequence]
\label{p12-83-c1-s8:cor:u006a-realizacion-sucesion}
Remark. For every sequence
\((b_k)_{k\geq1}\in\{-1,0,+1\}^{\mathbb N}\) there exists a cardinal history
whose observation satisfies
\[
 (a_3,a_6,a_9,\ldots)=(b_1,b_2,b_3,\ldots).
\]
\end{corollary}

\begin{proof}
Fixing a total order on \(\mathscr D_{\rm card}^3\), at stage \(k\) the
first element of \(\mathcal R_3(q_{k-1},b_k)\) is taken. The terminal state
of the block again belongs to \(\mathcal Q_{\rm loc}\), so that the
choice may be iterated indefinitely.
\end{proof}

The corollary proves that the sampled factor can be realized. It does not select the sections of \(\pi,e,\varphi\), fix the two intermediate trits or replace the compatibility of the enriched state.

\subsection{Cylinders of the \(2\) readers}

A ternary word of length \(L\), interpreted as the integer \(N\),
determines
\[
 I_3(N,L)=\left[\frac{N}{3^L},\frac{N+1}{3^L}\right).
\]
A prefix of \(K\) decimal triads, interpreted as the integer \(D\),
determines
\[
 I_{1000}(D,K)=
 \left[\frac{D}{1000^K},\frac{D+1}{1000^K}\right).
\]
We introduce the cross residues
\[
 A=N1000^K-D3^L,
 \qquad
 B=(N+1)1000^K-D3^L=A+1000^K.
\]

\begin{lemma}[Integer compatibility criterion]
\label{p12-83-c1-s8:lem:u006a-criterio-entero}
The cylinders \(I_3(N,L)\) and \(I_{1000}(D,K)\) intersect if and only if
\[
 \boxed{A<3^L,\qquad B>0.}
\]
Equivalently,
\[
 N1000^K<(D+1)3^L,
 \qquad
 (N+1)1000^K>D3^L.
\]
\end{lemma}

\begin{proof}
Two half-open intervals \([a,b)\) and \([c,d)\) intersect
exactly when \(a<d\) and \(c<b\). Multiplication by
\(3^L1000^K>0\) produces the two inequalities. Subtracting \(D3^L\) and
\(N1000^K\), respectively, gives the form in \(A,B\).
\end{proof}

\subsection{Completion operator \(C_{3,1000}\): ternary extension and preservation of the unit in \(10^3\) positions}

A new block of six trits is represented by
\[
 u\in\{0,\ldots,3^6-1\}=\{0,\ldots,728\}.
\]
The ternary update is
\[
 N'=729N+u,
 \qquad L'=L+6.
\]
Decimal publication has two cases. If the refined cylinder does not yet
determine a new triad, then \(K'=K\), \(D'=D\). If it determines
\(\delta\in\{0,\ldots,999\}\), then
\[
 K'=K+1,
 \qquad D'=1000D+\delta.
\]

\begin{definition}[Crossed state and exact transition]
\label{p12-83-c1-s8:def:u006a-c31000}
The cylinder state is the sextuple
\[
 \mathfrak c=(N,L,D,K,A,B)
\]
subject to
\[
 A=N1000^K-D3^L,
 \qquad B=A+1000^K.
\]
The transition \(\mathcal C_{3,1000}\) adjoins \(u\) and, when corresponds,
\(\delta\), and updates
\[
 (N,L,D,K,A,B)\longmapsto
 (N',L',D',K',A',B')
\]
through means of
\[
 A'=
 \begin{cases}
  729A+u1000^K,&K'=K,\\[1mm]
  729000A+u1000^{K+1}-\delta\,729\,3^L,&K'=K+1,
 \end{cases}
 \qquad
 B'=A'+1000^{K'}.
\]
The prolongation is admissible exactly when
\[
 A'<3^{L'},\qquad B'>0.
\]
\end{definition}

\begin{theorem}[Accuracy and Sufficiency of the Cross State]
\label{p12-83-c1-s8:thm:u006a-exactitud-c31000}
The transition of the \cref{p12-83-c1-s8:def:u006a-c31000} exactly preserves the
identities
\[
 A'=N'1000^{K'}-D'3^{L'},
 \qquad
 B'=(N'+1)1000^{K'}-D'3^{L'}.
\]
Consequently, \((A',B',L',K')\) suffices to decide the compatibility of the
new pair of cylinders, while \((N',D')\) preserves its reconstructive
provenance.
\end{theorem}

\begin{proof}
If a decimal triplet does not appear,
\[
 N'1000^K-D3^{L+6}
 =(729N+u)1000^K-729D3^L
 =729A+u1000^K.
\]
Si aparece \(\delta\),
\begin{align*}
 N'1000^{K+1}-D'3^{L+6}
 &=(729N+u)1000^{K+1}-(1000D+\delta)7293^L\\
 &=729000A+u1000^{K+1}-\delta\,729\,3^L.
\end{align*}
In both cases, increasing \(N'\) by one unit adds \(1000^{K'}\), which
yields \(B'\). The criterion of the \cref{p12-83-c1-s8:lem:u006a-criterio-entero} concludes.
\end{proof}

\begin{proposition}[Compatibility with concatenation]
\label{p12-83-c1-s8:prop:u006a-concatenacion-c31000}
Successively applying \(\mathcal C_{3,1000}\) to two blocks produces the same pair
\((N'',D'')\) and the same cross-residues as appending at once the two
ternary words and the decimal triads published in the same order.
\end{proposition}

\begin{proof}
The recurrences \(N\mapsto729N+u\) and
\(D\mapsto1000D+\delta\) are the laws of positional concatenation in
bases \(729\) and \(1000\). The
\cref{p12-83-c1-s8:thm:u006a-exactitud-c31000} identifies \(A,B\) with functions of those
four integers; hence its stagewise update coincides with the
evaluation subsequent to concatenation.
\end{proof}

\subsection{Decomposition of the Topological Phase Kernel into advance, rotation, carry, and memory operators}

The correct sequence is
\[
 (q,u)\xmapsto{\ R_3\ }b,
 \qquad
 (a_t)_{t\geq1}
 \xmapsto{\ \operatorname{Obs}_3\ }
 (a_{3k})_{k\geq1},
 \qquad
 (\mathfrak c,u,\delta)
 \xmapsto{\ \mathcal C_{3,1000}\ }
 \mathfrak c'.
\]
\(R_3\) proves local realizability; \(\operatorname{Obs}_3\) publishes a
subsequence; \(\mathcal C_{3,1000}\) transports compatibility between two
bases. None of the three can replace the survival relation
of the full TPK.

\begin{criterion}[Refutation criteria]
\label{p12-83-c1-s8:crit:u006a-refutacion}
The construction is refuted in its domain if any of the following
facts occurs:
\begin{enumerate}
 \item any of the \(486\) fibers \(\mathcal R_3(q,b)\) is empty or remains
       outside of the interval \([8,40]\);
 \item the null trit changes the active mode;
 \item \(R_3\) is identified with \(\operatorname{Obs}_3\);
 \item the recurrence of \(A'\) does not reproduce its direct definition;
 \item a declared compatible prolongation violates \(A'<3^{L'}\) or
       \(B'>0\);
 \item the decision depends on a numerical tolerance absent from the
       integer inequalities.
\end{enumerate}
\end{criterion}


% Unidad U006D. Registro dodecafásico integral y descomposición 1+5+6.

\section{Record dodecafasico integral, extension ciclica and decomposition \texorpdfstring{\(1+5+6\)}{1+5+6}}
\label{p12-83-c1-s9:chap:registro-dodecafasico-integral}

The observable period of one hundred eight positions contains twelve sectors of
nine positions. This equality of cardinalities does not suffice to describe the
temporal state: the ordered support, cyclic law, sector index, and
enriched register of transports, orientations, carries, and incidences must be
distinguished. This unit constructs the four objects and proves
the integral reconstruction of the twelve memory coordinates.

\subsection{Ordered support and cyclic coordinate}

For \(m\in\{1,\ldots,12\}\), define
\begin{equation}
 E_m:=\{9(m-1)+1,\ldots,9m\}.
 \label{p12-83-c1-s9:eq:u006d-bloques-nonadicos}
\end{equation}
The ordered support and the cyclic group are
\begin{equation}
 \Gamma_{108}:=
 E_1\sqcup\cdots\sqcup E_{12},
 \qquad
 \Theta_{108}:=\mathbb Z/108\mathbb Z.
 \label{p12-83-c1-s9:eq:u006d-gamma-theta}
\end{equation}
\(\Gamma_{108}\) preserves position, order and pertenencia to block;
\(\Theta_{108}\) preserves the law of translation. No are the same object.

For \(\theta\in\Theta_{108}\), let
\(\widetilde\theta\in\{0,\ldots,107\}\) be its representative. We define
\begin{equation}
 \beta(\theta):=
 \left[
 \left\lfloor\frac{\widetilde\theta}{9}\right\rfloor
 \right]_{12},
 \qquad
 r(\theta):=1+(\widetilde\theta\bmod9).
 \label{p12-83-c1-s9:eq:u006d-coordenadas-bloque-residuo}
\end{equation}
The first coordinate locates the duodecaphase sector; the second publishes
the positive position \(1,\ldots,9\) within the sector.

\subsection{Nonsplit cyclic extension}

Let \(C_n=\mathbb Z/n\mathbb Z\). The applications
\begin{equation}
 \iota:C_{12}\longrightarrow C_{108},
 \quad \iota([b]_{12})=[9b]_{108},
 \qquad
 p:C_{108}\longrightarrow C_9,
 \quad p([\theta]_{108})=[\theta]_9
 \label{p12-83-c1-s9:eq:u006d-mapas-extension}
\end{equation}
form the sequence
\begin{equation}
 0\longrightarrow C_{12}
 \xrightarrow{\;\iota\;}C_{108}
 \xrightarrow{\;p\;}C_9
 \longrightarrow0.
 \label{p12-83-c1-s9:eq:u006d-extension}
\end{equation}
For the set-theoretic section
\[
 s([r]_9):=[\widetilde r]_{108},
 \qquad 0\leq\widetilde r<9,
\]
the defect of additivity is
\begin{equation}
 c(r,r'):=
 \left[
 \left\lfloor
 \frac{\widetilde r+\widetilde r'}9
 \right\rfloor
 \right]_{12},
 \qquad
 s(r)+s(r')
 =s(r+r')+\iota(c(r,r')).
 \label{p12-83-c1-s9:eq:u006d-cociclo-extension}
\end{equation}

\begin{theorem}[Accuracy, non-division, and cohomological class]
\label{p12-83-c1-s9:thm:u006d-extension-no-escindida}
The sucesion \eqref{p12-83-c1-s9:eq:u006d-extension} is exacta and not it escinde. For the
action trivial,
\begin{equation}
 H^2(C_9,C_{12})
 \simeq C_{12}/9C_{12}
 \simeq C_3,
 \label{p12-83-c1-s9:eq:u006d-cohomologia-extension}
\end{equation}
and the cocycle \eqref{p12-83-c1-s9:eq:u006d-cociclo-extension} represents its generating
class.
\end{theorem}

\begin{proof}
The image of \(\iota\) is the subgroup of multiples of nine, which coincides
with \(\ker p\); \(p\) is surjective. This proves exactness. If the
extension were split, \(C_{108}\) would be isomorphic to
\(C_{12}\times C_9\), but
\[
 \exp(C_{108})=108,
 \qquad
 \exp(C_{12}\times C_9)=\operatorname{mcm}(12,9)=36.
\]
The cocycle identity is the Euclidean division of \(\widetilde r+\widetilde r'\) by nine.
For a cyclic group with trivial action, \(H^2(C_9,C_{12})\simeq C_{12}/9C_{12}\); the unit carry
projects to \([1]\), which generates the quotient of order three.
\end{proof}

The non-splitting typifies the return:
\[
 g_0(\theta+9)=g_0(\theta),
 \qquad
 \beta(\theta+9)=\beta(\theta)+1\pmod{12}.
\]
Nine iterations close the visible phase and shift one sector; one hundred
eight close both coordinates of the clock. The enriched state further preserves
the vertical memory, cylinder, boundary, and genealogy.

\subsection{Enhanced temporal registration: joint encoding of phase, sheet, boundary, and memory}

\begin{definition}[Oriented dodecaphonic register]
\label{p12-83-c1-s9:def:u006d-registro-dodecafasico}
For an enriched history, let \(\tau_m\) be the ordered subroute on
\(E_m\), \(\sigma_m\) its orientation, \(\kappa_m\) the carry crossing the
block boundary, and \(\Lambda_m\) its local incidence ledger. The
complete record is
\begin{equation}
 \mathcal R_{12}:=
 \bigl(
 (E_m,\tau_m,\sigma_m,\kappa_m,\Lambda_m)
 \bigr)_{m=1}^{12}.
 \label{p12-83-c1-s9:eq:u006d-registro-r12}
\end{equation}
\end{definition}

Thus four levels are distinguished:
\begin{equation}
 \boxed{
 \mathcal R_{12}
 \neq C_{12}
 \neq\Theta_{108}
 \neq\Gamma_{108}.}
 \label{p12-83-c1-s9:eq:u006d-cuatro-objetos}
\end{equation}
The first is an enriched register; the second, a sector index; the
third, a group of translations; the fourth, an ordered support. There exist
projections between them, but no identifications.

Let \(q_{\rm mem}\in\mathbb Z_{\geq0}\) be the unbounded number of nonadic
turns. The duodecaphase coordinate is only its residue
\[
 \beta=[q_{\rm mem}]_{12}.
\]
After one hundred and eight steps,
\[
 \beta\longmapsto\beta,
 \qquad
 q_{\rm mem}\longmapsto q_{\rm mem}+12.
\]
The calendar returns; the vertical memory is not reset.

\subsection{Signed event extractor}

Let \((d_t)_{t=1}^{108}\) be the digital record on
\(\Gamma_{108}\). We define
\begin{equation}
 \mathcal D_{\rm evt}:=\{2,4,6,8\},
 \qquad
 \Sigma_{12}:=(+,+,+,-,-,-,+,+,+,-,-,-).
 \label{p12-83-c1-s9:eq:u006d-datos-eventos}
\end{equation}
For \(m=0,\ldots,11\), let
\begin{equation}
 e_m:=
 \Sigma_{12}(m+1)\,
 \#\{t\in E_{m+1}:d_t\in\mathcal D_{\rm evt}\}
 \in\mathbb Z.
 \label{p12-83-c1-s9:eq:u006d-eventos-firmados}
\end{equation}
The application preserves the local census and the sectorial orientation as separate factors.

We pair the two semicrowns by means of
\begin{equation}
 p_i:=e_i+e_{i+6},
 \qquad
 a_i:=e_i-e_{i+6},
 \qquad 0\leq i\leq5.
 \label{p12-83-c1-s9:eq:u006d-pares-antipodales}
\end{equation}
We define the total census, the five differences with respect to the sixth pair,
and the six antisymmetries:
\begin{equation}
 s_C:=\sum_{i=0}^{5}p_i,
 \qquad
 M_i:=p_i-p_5\quad(0\leq i\leq4),
 \qquad
 A_6:=(a_0,\ldots,a_5).
 \label{p12-83-c1-s9:eq:u006d-coordenadas-uno-cinco-seis}
\end{equation}

\begin{definition}[Integral reader \(1+5+6\)]
\label{p12-83-c1-s9:def:u006d-lector-integral}
The linear reader is
\begin{equation}
 \mathcal L_{12}:\mathbb Z^{12}\longrightarrow\mathbb Z^{12},
 \qquad
 \mathcal L_{12}(e)
 :=
 (s_C,M_0,\ldots,M_4,a_0,\ldots,a_5).
 \label{p12-83-c1-s9:eq:u006d-lector-l12}
\end{equation}
The coordinate distribution is
\begin{equation}
 \underbrace{1}_{s_C}
 +\underbrace{5}_{M_0,\ldots,M_4}
 +\underbrace{6}_{a_0,\ldots,a_5}
 =12.
 \label{p12-83-c1-s9:eq:u006d-uno-cinco-seis}
\end{equation}
\end{definition}

\begin{theorem}[Exact reconstruction and integral image]
\label{p12-83-c1-s9:thm:u006d-inversa-l12}
The reader \(\mathcal L_{12}\) is injective. On its image,
\begin{align}
 p_5&=
 \frac{s_C-\sum_{i=0}^{4}M_i}{6},
 &
 p_i&=p_5+M_i\quad(0\leq i\leq4),
 \label{p12-83-c1-s9:eq:u006d-inversa-p}\\
 e_i&=\frac{p_i+a_i}{2},
 &
 e_{i+6}&=\frac{p_i-a_i}{2}
 \quad(0\leq i\leq5).
 \label{p12-83-c1-s9:eq:u006d-inversa-e}
\end{align}
An output tuple belongs to \(\operatorname{im}\mathcal L_{12}\) if and
only if
\begin{equation}
 s_C-\sum_{i=0}^{4}M_i\equiv0\pmod6
 \label{p12-83-c1-s9:eq:u006d-condicion-divisibilidad}
\end{equation}
and, for the reconstructed \(p_i\),
\begin{equation}
 p_i\equiv a_i\pmod2
 \qquad(0\leq i\leq5).
 \label{p12-83-c1-s9:eq:u006d-condicion-paridad}
\end{equation}
\end{theorem}

\begin{proof}
From \(M_i=p_i-p_5\) is obtained
\[
 s_C=\sum_{i=0}^{4}(p_5+M_i)+p_5
 =6p_5+\sum_{i=0}^{4}M_i,
\]
which gives the first equation of \eqref{p12-83-c1-s9:eq:u006d-inversa-p}; the others are
immediate. Adding and subtracting
\(p_i=e_i+e_{i+6}\) and \(a_i=e_i-e_{i+6}\) produces
\eqref{p12-83-c1-s9:eq:u006d-inversa-e}. Divisibility by six and the six
parity congruences are precisely the necessary and
sufficient conditions for all reconstructed coordinates to be integers.
\end{proof}

Division by six appears after the oriented census, not during it. The
integrality conditions for this unit are exactly
\eqref{p12-83-c1-s9:eq:u006d-condicion-divisibilidad} and
\eqref{p12-83-c1-s9:eq:u006d-condicion-paridad}; no Smith normal form is attributed here to a
matrix that has not previously been defined.

\subsection{\(2\) dodecaphasic words of different provenance}

The distinction must be preserved.
\begin{equation}
 U_{12}^{\rm cat}:=U_6\Vert U_6,
 \qquad
 U_{12}^{\rm sgn}:=\mathcal H_{12}(K).
 \label{p12-83-c1-s9:eq:u006d-dos-palabras-doce}
\end{equation}
The first concatenates two catalog emissions of length six. The second
is a signed coordinate reversibly linked to \(K\) by the operator
\(\mathcal H_{12}\). Their common length does not imply the same provenance,
semantics, or transport law.

\begin{criterion}[Falsification criteria]
\label{p12-83-c1-s9:crit:u006d-refutacion}
The construction is refuted if the sequence
\eqref{p12-83-c1-s9:eq:u006d-extension} is not exact, if its cohomology class is
trivial, if a tuple in the image violates
\eqref{p12-83-c1-s9:eq:u006d-condicion-divisibilidad} or
\eqref{p12-83-c1-s9:eq:u006d-condicion-paridad}, if the inverse formulas do not recover
the twelve initial integers, or if the return of \(\beta\) is interpreted as a
reset of \(q_{\rm mem}\).
\end{criterion}


% Unidad U006B. Funtor ciclotomico y elevacion modular de caminos.

\section{Cyclotomic phase functor--length and modular elevation of paths}
\label{p12-83-c1-s10:chap:funtor-ciclotomico-elevacion-caminos}

The odometric quotient \(\mathcal O_{9,12}\) simultaneously preserves the nonadic position, the dodecaphase register, and the integral length of each arrow. In this unit, its phase is realized over fibers \(A_2\), and the cardinal lift that replaces each edge with four oriented edges is proved. The cyclotomic representation and the absolute length will be retained as distinct coordinates.

\subsection{Coxeter representation of order \(3\) on the dodecaphasic fibers of type \(A_2\)}

Let
\begin{equation}
 G_{A_2}:=
 \begin{pmatrix}2&-1\\-1&2\end{pmatrix},
 \qquad
 C:=
 \begin{pmatrix}0&-1\\1&-1\end{pmatrix}.
 \label{p12-83-c1-s10:eq:u006b-dato-coxeter}
\end{equation}
By the \cref{p12-83-c1-s5:prop:coxeter-a2-autonomo},
\begin{equation}
 C^3=I_2,
 \qquad
 C^{\mathsf T}G_{A_2}C=G_{A_2},
 \qquad
 \det(XI_2-C)=X^2+X+1.
 \label{p12-83-c1-s10:eq:u006b-identidades-coxeter}
\end{equation}
Thus, \(C\) realizes the character cyclotomic of \(C_3\) on \(A_2\).

For each \(b\in\mathbb Z/12\mathbb Z\), let \(A_{2,b}\) be a copy of
\(A_2\), let
\[
 \iota_b:A_{2,b}\xrightarrow{\;\sim\;}A_2
\]
an integral identification, and let
\(P_b\in O(G_{A_2})\) be a frame that depends on \(b\), but not on the residue
\(r\). We define
\begin{equation}
 \rho_b:=
 \iota_b^{-1}P_bCP_b^{-1}\iota_b,
 \qquad
 \rho_b^3=\operatorname{id}_{A_{2,b}}.
 \label{p12-83-c1-s10:eq:u006b-rho-fibra}
\end{equation}
The independencia respect of \(r\) liga the character to the fiber
dodecaphasic; \(r\) preserves the position internal of the turn nonadic.

\subsection{Odometer category of the TPK: cylinders, truncations, and sequence morphisms}

Recall that \(\mathcal O_{9,12}\) has objects
\[
 (b,r)\in\mathbb Z/12\mathbb Z\times\{0,\ldots,8\}.
\]
For \(n\in\mathbb N\), the flecha unique of length \(n\) is
\begin{equation}
 \gamma_{(b,r),n}:
 (b,r)\longrightarrow
 \left(
 b+\kappa_r(n),[r+n]_9
 \right),
 \qquad
 \kappa_r(n):=
 \left\lfloor\frac{r+n}{9}\right\rfloor.
 \label{p12-83-c1-s10:eq:u006b-flecha-odometro}
\end{equation}
The composition relies on
\begin{equation}
 \kappa_r(n+m)
 =\kappa_r(n)+\kappa_{[r+n]_9}(m).
 \label{p12-83-c1-s10:eq:u006b-cociclo-odometro}
\end{equation}
The length \(\ell(\gamma_{(b,r),n})=n\) is additive.

\begin{definition}[Cyclotomic Phase--Length Functor]
\label{p12-83-c1-s10:def:u006b-funtor-ciclotomico}
Let \(\mathbf{Rep}_{\mathbb C}(C_3)\) be the category of complex
representations of \(C_3\), and let \(\mathbf B\mathbb N\) be the monoidal category
with a single object whose endomorphisms are \((\mathbb N,+)\). We define
\begin{equation}
 \mathfrak F_{\rm cyc}:
 \mathcal O_{9,12}
 \longrightarrow
 \mathbf{Rep}_{\mathbb C}(C_3)\times\mathbf B\mathbb N
 \label{p12-83-c1-s10:eq:u006b-signatura-funtor}
\end{equation}
by
\[
 \mathfrak F_{\rm cyc}(b,r)
 =(A_{2,b}\otimes_{\mathbb Z}\mathbb C,\rho_b)
\]
and
\begin{equation}
 \mathfrak F_{\rm cyc}(\gamma)
 :=
 \left(
 \iota_{b'}^{-1}P_{b'}C^{\ell(\gamma)}
 P_b^{-1}\iota_b,
 \ell(\gamma)
 \right)
 \label{p12-83-c1-s10:eq:u006b-flecha-funtor}
\end{equation}
for \(\gamma:(b,r)\to(b',r')\).
\end{definition}

\begin{theorem}[Functoriality and separation of the records]
\label{p12-83-c1-s10:thm:u006b-functorialidad}
The assignment of the \cref{p12-83-c1-s10:def:u006b-funtor-ciclotomico} is a functor. Its
first component preserves only the length class modulo three; the
second preserves the complete integer and, by reduction, determines that class.
The cyclotomic component does not allow the absolute length to be reconstructed.
\end{theorem}

\begin{proof}
If \(\gamma_1:(b_0,r_0)\to(b_1,r_1)\) and
\(\gamma_2:(b_1,r_1)\to(b_2,r_2)\) have lengths \(n_1,n_2\), the
intermediate trivialization cancels:
\begin{align}
 &\bigl(
 \iota_{b_2}^{-1}P_{b_2}C^{n_2}P_{b_1}^{-1}\iota_{b_1}
 \bigr)
 \bigl(
 \iota_{b_1}^{-1}P_{b_1}C^{n_1}P_{b_0}^{-1}\iota_{b_0}
 \bigr)\notag\\
 &\qquad=
 \iota_{b_2}^{-1}P_{b_2}C^{n_1+n_2}
 P_{b_0}^{-1}\iota_{b_0}.
 \label{p12-83-c1-s10:eq:u006b-composicion-funtor}
\end{align}
The second component composes by addition. Since \(C^n\) depends only on
\([n]_3\), whereas \(\ell=n\) preserves the integer, the two coordinates are not
identified.
\end{proof}

If
\[
 Q:\mathsf{TPK}\longrightarrow\mathcal O_{9,12}
\]
is the odometer functor of U006, the typed realization is
\begin{equation}
 \mathfrak F_{\rm cyc}\circ Q:
 \mathsf{TPK}\longrightarrow
 \mathbf{Rep}_{\mathbb C}(C_3)\times\mathbf B\mathbb N.
 \label{p12-83-c1-s10:eq:u006b-compuesto-tpk-ciclotomico}
\end{equation}

\subsection{Modular elevation of cardinal paths}

Let \(\mathcal A_{N,d}\) be the category libre of paths cardinal of
\((\mathbb Z/N\mathbb Z)^d\). Let
\[
 N\ge2,\quad u\ge1,\quad\gcd(u,N)=1,
 \quad k:=\operatorname{ord}_N(u),
 \quad q:=\frac{u^k-1}{N}.
\]

\begin{definition}[Cardinal Substitution by \(u\)]
\label{p12-83-c1-s10:def:u006b-sustitucion-cardinal}
The functor
\[
 \mathcal S_u:\mathcal A_{N,d}\longrightarrow\mathcal A_{N,d}
\]
sends an object \(p\) to \(up\) and each oriented edge to the ordered concatenation of \(u\) copies with the same direction.
\end{definition}

\begin{theorem}[Modular lifting to endomorphism of paths]
\label{p12-83-c1-s10:thm:u006b-elevacion-modular}
For every object \(p\), path \(\gamma\), and
\[
 \kappa_{N,r}(n):=
 \left\lfloor\frac{r+n}{N}\right\rfloor,
\]
we have
\begin{align}
 \mathcal S_u^k(p)&=p,
 \label{p12-83-c1-s10:eq:u006b-retorno-objetos}\\
 \ell(\mathcal S_u^k\gamma)&=u^k\ell(\gamma),
 \label{p12-83-c1-s10:eq:u006b-longitud-sustituida}\\
 \kappa_{N,r}(u^kn)-\kappa_{N,r}(n)&=qn.
 \label{p12-83-c1-s10:eq:u006b-diferencia-vueltas}
\end{align}
\end{theorem}

\begin{proof}
The category is free, so the images of objects and generators
determine a unique functor. Each map multiplies the length by
\(u\). Since \(u^k=1+qN\), the objects return modulo \(N\), and
\[
 \left\lfloor
 \frac{r+n+qNn}{N}
 \right\rfloor
 =
 \left\lfloor\frac{r+n}{N}\right\rfloor+qn.
\]
\end{proof}

In the nonadic case,
\begin{equation}
 (N,u,k,q)=(9,4,3,7),
 \qquad
 4^3=64=1+7\cdot9.
 \label{p12-83-c1-s10:eq:u006b-parametros-nonadicos}
\end{equation}
Therefore,
\begin{equation}
 \boxed{
 \mathcal S_4^3(e)
 \text{ preserves the endpoints and contains seven nonadic turns
 oriented per initial edge}.}
 \label{p12-83-c1-s10:eq:u006b-siete-vueltas}
\end{equation}
The endpoints and the residual position return under three maps; the
enriched state need not return. Without a proof of
invertibility, \(\mathcal S_4^3\) is a vertical endomorphism of paths,
not a monodromy.

\begin{remark}[Separation with respect to clock dilation]
\label{p12-83-c1-s10:rem:u006b-s4-no-d4}
\(\mathcal S_4\) replaces objects and edges in a path category.
The map \(D_4\) of U006C will multiply the coordinate
\(\Theta_{108}\) by four. They share an integer factor, but not a domain, codomain, or
composition law.
\end{remark}

\begin{criterion}[Criteria of refutation]
\label{p12-83-c1-s10:crit:u006b-refutacion}
The block remains refuted if fails any of the identities
\eqref{p12-83-c1-s10:eq:u006b-identidades-coxeter},
\eqref{p12-83-c1-s10:eq:u006b-cociclo-odometro},
\eqref{p12-83-c1-s10:eq:u006b-composicion-funtor} or
\eqref{p12-83-c1-s10:eq:u006b-parametros-nonadicos}; also remains badly typed if it
identify the phase \(C^n\), the length \(n\), the substitution
\(\mathcal S_4\) and the dilation \(D_4\).
\end{criterion}
}%


% Unidad U004. Redacción autónoma desde los generadors TRIT de 1063/live.

\section{TRIT as the minimal unit of topological information: ternary orientation, carry memory, and temporal order}
\label{p12-83-c1-s11:chap:trit-levantado-regimenes}

The mixed curvature of APP produced three oriented values: zero on the pure sheets and opposite signs when the sum--product order is exchanged. The TRIT types that triad and preserves, in addition to the visible residue, the quotient required to reconstruct the lifted integer. This unit explicitly separates four levels: ternary class, balanced representative, carry, and quadratic regime.

\begin{definition}[Topological Information Unit TRIT]
The \emph{TRIT} is the elementary oriented unit whose state space is
\(\mathsf{TRIT}=\{-1,0,+1\}\). Its native logarithmic capacity is
\(\log 3\), and its expression in binary units is \(\log_2 3\). Unlike
a bit, its three states do not reduce to a cardinal encoding: they respectively type
opening, threshold, and return, and preserve the carry quotient
required to reconstruct the transition that produced the visible
residue.
\end{definition}

\subsection{Oriented ternary system and balanced representation of TRIT states}

We define the oriented set
\[
 \mathsf{TRIT}:=\{-1,0,+1\}.
\]
The balanced chart is the bijection of sets.
\[
 \operatorname{bal}:\mathbb F_3\longrightarrow\mathsf{TRIT},
 \qquad
 [0]_3\mapsto0,\qquad[1]_3\mapsto+1,\qquad[2]_3\mapsto-1.
\]
It is not of an equality of types. \(\mathbb F_3\) carries sum and product
of field; \(\mathsf{TRIT}\) carries order and orientation. Every operation
transported between both must indicate the chart used.

The orientation of the nine phases is
\[
 \varepsilon_9(r)=
 \begin{cases}
 +1,&r\in\{1,4,7\},\\
 -1,&r\in\{2,5,8\},\\
 0,&r\in\{3,6,9\}.
 \end{cases}
\]
His complete table is
\[
\begin{array}{c|ccccccccc}
r&1&2&3&4&5&6&7&8&9\\\hline
\varepsilon_9(r)&+1&-1&0&+1&-1&0&+1&-1&0
\end{array}.
\]
Periodicity of three does not eliminate the nonadic phase: the phases
\(1,4,7\), for example, have the same orientation and different positions
in the circuit.

\subsection{Balanced ternary division algorithm and uniqueness of the quotient–residue pair}

\begin{theorem}[Local ternary lifting]
\label{p12-83-c1-s11:thm:trit-levantamiento-local}
For each \(n\in\{-4,-3,\ldots,4\}\) there exists a single pair
\((r,c)\in\mathsf{TRIT}^2\) such that
\[
 n=r+3c.
\]
\end{theorem}

\begin{proof}
The existence is verified in the table.
\[
\begin{array}{c|rrrrrrrrr}
n&-4&-3&-2&-1&0&1&2&3&4\\\hline
r&-1&0&+1&-1&0&+1&-1&0&+1\\
c&-1&-1&-1&0&0&0&+1&+1&+1
\end{array}.
\]
If \(r+3c=r'+3c'\), then \(r-r'=3(c'-c)\). Since
\(r-r'\in\{-2,-1,0,1,2\}\), the only possible multiple of three is zero;
therefore \(r=r'\) and \(c=c'\).
\end{proof}

The visible projection \((r,c)\mapsto r\) has three local states over the value zero:
\[
 0^-:=(0,-1),
 \qquad
 0^0:=(0,0),
 \qquad
 0^+:=(0,+1).
\]
They publish the same residue, but reconstruct \(-3,0,+3\). This triple fiber
is the minimal form of a distinction that will reappear in the enriched state
of the Topological Phase Kernel.

\begin{corollary}[Balanced Iteration]
\label{p12-83-c1-s11:cor:trit-iteracion-balanceada}
Every integer \(N\) admits a finite balanced ternary expansion
\[
 N=r_0+3r_1+\cdots+3^mr_m,
 \qquad r_j\in\mathsf{TRIT}.
\]
It is obtained by iterating the division \(N_k=r_k+3N_{k+1}\), where \(r_k\) is the
balanced representative of \(N_k\bmod3\).
\end{corollary}

\begin{proof}
Euclidean division modulo three yields a residue in \(\{0,1,2\}\); replacing \(2\) with \(-1\) increases the quotient by one and produces \(r_k\in\mathsf{TRIT}\). The absolute value of the quotient decreases outside a finite set, so the process terminates. Uniqueness follows by comparing the first differing digit and reducing modulo three.
\end{proof}

\subsection{Composition and carry cocycle}

For \(r,s\in\mathsf{TRIT}\), we define \(r\oplus s\in\mathsf{TRIT}\) and
\(\kappa_3(r,s)\in\mathbb Z\) by the unique equality
\[
 r+s=(r\oplus s)+3\kappa_3(r,s).
\]
The table is
\[
\begin{array}{c|ccc}
\oplus&-1&0&+1\\\hline
-1&+1&-1&0\\
0&-1&0&+1\\
+1&0&+1&-1
\end{array},
\qquad
\begin{array}{c|ccc}
\kappa_3&-1&0&+1\\\hline
-1&-1&0&0\\
0&0&0&0\\
+1&0&0&+1
\end{array}.
\]

\begin{proposition}[Ternary cocyclic identity]
\label{p12-83-c1-s11:prop:trit-cociclo}
For \(r,s,t\in\mathsf{TRIT}\),
\[
 \kappa_3(r,s)+\kappa_3(r\oplus s,t)
 =\kappa_3(s,t)+\kappa_3(r,s\oplus t).
\]
\end{proposition}

\begin{proof}
Both sides are the total quotient that appears when writing
\(r+s+t\) as a balanced representative plus a multiple of three. The
uniqueness of the lift proves the equality.
\end{proof}

If \((r,c)\) and \((s,d)\) are states levantados, its sum is
\[
 (r,c)\boxplus(s,d)
 :=\bigl(r\oplus s,\ c+d+\kappa_3(r,s)\bigr).
\]
The \cref{p12-83-c1-s11:prop:trit-cociclo} proves that \(\boxplus\) is associative and that the
reconstruction \(\operatorname{rec}_3(r,c)=r+3c\) satisfies
\[
 \operatorname{rec}_3(x\boxplus y)
 =\operatorname{rec}_3(x)+\operatorname{rec}_3(y).
\]

\subsection{Orientation and Temporal Order Involution}

For to tritic word \(\tau=(\tau_1,\ldots,\tau_n)\), define
\[
 \mathcal C_{\rm rev}(\tau_1,\ldots,\tau_n)
 :=(-\tau_n,\ldots,-\tau_1).
\]

\begin{lemma}[Oriented reversal]
\label{p12-83-c1-s11:lem:trit-reversion}
The map \(\mathcal C_{\rm rev}\) is an involution. It interchanges the
sectors \(+1\) and \(-1\) and fixes the sector \(0\).
\end{lemma}

\begin{proof}
Applying it twice reverses the order twice and changes every sign twice.
The initial word is thereby recovered. The remaining assertions are
immediate in each component.
\end{proof}

When a time-oriented parameter is fixed, the TPK convention is
\[
 +1:\text{return with memory},
 \qquad
 0:\text{present threshold},
 \qquad
 -1:\text{prolongation opening}.
\]
This assignment does not add a fourth temporal state. The direction is preserved
in the sign and the exact position is preserved in the nonadic phase.

\subsection{Universal family of quadratic algebras}

For \(\tau\in\mathsf{TRIT}\), define
\[
 \mathbb A_\tau:=\mathbb R[J_\tau]/(J_\tau^2+\tau I).
\]
The three cases are:
\[
\begin{array}{c|c|c|c}
\tau&J_\tau^2&\text{regime}&\exp(tJ_\tau)\\\hline
+1&-I&\text{elliptic}&\cos t\,I+\sin t\,J_{+1}\\
0&0&\text{parabolic}&I+tJ_0\\
-1&+I&\text{hyperbolic}&\cosh t\,I+\sinh t\,J_{-1}
\end{array}.
\]

\begin{theorem}[Uniform quadratic exponential]
\label{p12-83-c1-s11:thm:trit-exponencial-uniforme}
If \(J_\tau^2=-\tau I\), then
\[
 \exp(tJ_\tau)
 =C_\tau(t)I+S_\tau(t)J_\tau,
\]
where
\[
\begin{array}{c|ccc}
\tau&+1&0&-1\\\hline
C_\tau(t)&\cos t&1&\cosh t\\
S_\tau(t)&\sin t&t&\sinh t
\end{array}.
\]
\end{theorem}

\begin{proof}
We separate the exponential series into even and odd powers. The relations
\(J_{+1}^{2m}=(-1)^mI\), \(J_0^2=0\), and
\(J_{-1}^{2m}=I\) respectively produce the cosine and sine series,
the linear truncation, and the hyperbolic series.
\end{proof}

On \(\mathbb F_3\) the analogues appear
\[
 \mathbb F_9=\mathbb F_3[\iota]/(\iota^2+1),
 \qquad
 \mathbb D_3=\mathbb F_3[\epsilon]/(\epsilon^2),
 \qquad
 \mathbb S_3=\mathbb F_3[j]/(j^2-1).
\]
The first is a field of nine elements; the second contains a nonzero
nilpotent; the third decomposes through the idempotents
\((1\pm j)/2\) and is isomorphic to \(\mathbb F_3\times\mathbb F_3\). The
coincidence of cardinality does not identify these three algebras.

The integral source of the elliptic regime is
\[
 \mathbb Q_{\mathbb Z}:=\mathbb Z[u]/(u^2+1).
\]
Reduction modulo three sends \(u\) to \(\iota\) and produces
\(\mathbb F_9\). The extension of scalars to \(\mathbb R\) and a
\(u\mapsto J_{+1}\) representation produce the real elliptic realizer.
There is therefore no untyped map \(\mathbb F_9\to\mathbb R\): both
realizations follow from a common integral source.

\subsection{Matrix representations of the TRIT selected by the sign of \(J_\tau^2\)}

Let
\[
 C=\begin{pmatrix}0&-1\\1&-1\end{pmatrix},
 \qquad C^3=I,
\]
and let us define
\[
 J_e:=\frac{2C+I}{\sqrt3}
 =\frac1{\sqrt3}\begin{pmatrix}1&-2\\2&-1\end{pmatrix}.
\]
A direct multiplication gives
\[
 (2C+I)^2=-3I,
 \qquad J_e^2=-I.
\]
The parabolic realizer is
\[
 J_p:=N=\begin{pmatrix}0&1\\0&0\end{pmatrix},
 \qquad N^2=0.
\]
For the hyperbolic regime we take
\[
 F_{\rm av}=\begin{pmatrix}0&1\\1&1\end{pmatrix},
 \qquad
 A:=F_{\rm av}^2=\begin{pmatrix}1&1\\1&2\end{pmatrix},
\]
and
\[
 J_h:=\frac{2A-3I}{\sqrt5}
 =\frac1{\sqrt5}\begin{pmatrix}-1&2\\2&1\end{pmatrix}.
\]
Then
\[
 (2A-3I)^2=5I,
 \qquad J_h^2=I.
\]

The bilinear form
\[
 G=\begin{pmatrix}2&1\\1&-2\end{pmatrix}
\]
satisfies
\[
 A^TGA=G.
\]
Moreover,
\[
 \cosh(2\log\varphi)=\frac32,
 \qquad
 \sinh(2\log\varphi)=\frac{\sqrt5}{2},
\]
for which
\[
 A=\exp(2\log\varphi\,J_h)
  =\cosh(2\log\varphi)I+\sinh(2\log\varphi)J_h.
\]
This identity will be recovered causally in the biography of
\(\varphi\); here only the matrix representation is verified.

The three exponential controls are
\[
 \exp(\pi J_e)+I=0,
 \qquad
 \exp(J_p)=I+J_p,
 \qquad
 \exp(2\log\varphi\,J_h)=F_{\rm av}^2.
\]
The presence of \(\pi\) and \(\varphi\) in these verification identities
does not authorize using them as constructors of those constants. Their
constructors will appear later from the nonadic genealogy.

\subsection{Local coupling APP--TRIT}

With the exchange
\[
 R=\begin{pmatrix}0&1\\1&0\end{pmatrix}
\]
we define, for \(\tau\in\mathsf{TRIT}\),
\[
 B_\tau=(4+3\tau)I+(5-3\tau)R.
\]
Since \(P_\pm=(I\pm R)/2\),
\[
 B_\tau=9P_++(6\tau-1)P_-.
\]

\begin{theorem}[Spectral Recovery of Orientation]
\label{p12-83-c1-s11:thm:app-trit-recuperacion-espectral}
The spectrum of \(B_\tau\) is
\[
 \operatorname{Spec}(B_\tau)=\{9,6\tau-1\}.
\]
The differential eigenvalue \(\lambda_-\) uniquely determines the TRIT:
\[
 \tau=\frac{\lambda_-+1}{6}.
\]
Moreover, the block central of APP is \(B_{+1}\).
\end{theorem}

\begin{proof}
The projectors \(P_+\) and \(P_-\) diagonalize \(R\) with eigenvalues
\(+1\) and \(-1\). The spectral formula is obtained by substitution. The second
identity solves for \(\tau\). For \(\tau=+1\),
\(B_{+1}=7I+2R=\begin{psmallmatrix}7&2\\2&7\end{psmallmatrix}\), which is
the unique block of \cref{p12-83-c1-s1:thm:app-bloque-central}.
\end{proof}

Therefore, the TRIT not it anade as retrospective label: can be recovered
of the differential mode of a family constructed on the internal exchange
of APP.

\subsection{Obstructions of TRIT: residue--representative--carry separation, involutivity, and spectral recovery of \(\tau\)}

\begin{criterion}[Refutation Criteria for the TRIT]
The unit is refuted if:
\begin{enumerate}
 \item it identify class residual, representative balanceado and carry;
 \item it is asserted that \(0^-\), \(0^0\), and \(0^+\) are the same state;
 \item the composition \(\boxplus\) ceases to reconstruct the integer sum;
 \item \(\mathcal C_{\rm rev}\) is not involutive;
 \item the three algebras are identified because they have the same cardinal;
 \item some realizer fails to satisfy \(J_\tau^2=-\tau I\);
 \item the value of \(\tau\) is chosen after observing the outcome and is not
       recovered from the spectrum of \(B_\tau\).
\end{enumerate}
\end{criterion}

APP provides the support, the multiplicative matrix, and the internal additive
law; TRIT provides orientation,
lift, and quadratic type. The next unit will construct the
Topological Phase Kernel as the transport category of these data, without
confusing the internal selector with the complete state.

\HMTDeferredTPKBase
\HMTDeferredTPKDevelopment


% Unidad U006C. Dilatación aritmética, cociclo de acarreo y
% semiconjugación cuaternaria.

\section{Quaternary dilation of the register, observable obstruction, and refined semiconjugacy}
\label{p12-83-c1-s12:chap:dilatacion-cuaternaria-semiconjugacion}

The cardinal substitution \(\mathcal S_4\) constructed in the preceding unit
acts on paths. In this unit, a second operation, arithmetic in
nature, is defined on the register
\(\mathbb Z/12\mathbb Z\times\mathbb Z/9\mathbb Z\), and a third,
categorical operation subdivides each observable transition into four
microtransitions. The aim is not to merge them, but to prove the exact
relationship between their realizations.

\subsection{Dilation on the product register}

Let
\[
 \mathcal Q_{12,9}:=
 \mathbb Z/12\mathbb Z\times\mathbb Z/9\mathbb Z.
\]
For \(r\in\mathbb Z/9\mathbb Z\), we denote its normalized representative by
\(\widetilde r\in\{0,\ldots,8\}\).

\begin{definition}[Quaternary dilation of the register]
\label{p12-83-c1-s12:def:u006c-dilatacion-registro}
We define
\begin{equation}
 D_4(b,r):=
 \left(
 \left[
  4b+\left\lfloor\frac{4\widetilde r}{9}\right\rfloor
 \right]_{12},
 [4r]_9
 \right).
 \label{p12-83-c1-s12:eq:u006c-d4-producto}
\end{equation}
The total coordinate is
\begin{equation}
 \Theta:\mathcal Q_{12,9}\xrightarrow{\;\sim\;}
 \mathbb Z/108\mathbb Z,
 \qquad
 \Theta(b,r):=[9b+\widetilde r]_{108}.
 \label{p12-83-c1-s12:eq:u006c-coordenada-total}
\end{equation}
\end{definition}

\begin{proposition}[Realization by modular multiplication]
\label{p12-83-c1-s12:prop:u006c-d4-total}
In the total coordinate,
\begin{equation}
 \boxed{
 \Theta\circ D_4\circ\Theta^{-1}(\theta)=[4\theta]_{108}.}
 \label{p12-83-c1-s12:eq:u006c-d4-total}
\end{equation}
In particular,
\begin{equation}
 \operatorname{im}(D_4)
 =4\mathbb Z/108\mathbb Z
 \simeq\mathbb Z/27\mathbb Z,
 \label{p12-83-c1-s12:eq:u006c-imagen-d4}
\end{equation}
and the induced restriction on the image has order nine.
\end{proposition}

\begin{proof}
For \(\theta=9b+\widetilde r\),
\[
 4\theta
 =9\left(
  4b+\left\lfloor\frac{4\widetilde r}{9}\right\rfloor
 \right)
 +(4\widetilde r\bmod9).
\]
This is exactly the quotient--residue decomposition of
\eqref{p12-83-c1-s12:eq:u006c-d4-producto}. The image consists of the classes that are
multiples of four. Dividing by four yields
\(\mathbb Z/27\mathbb Z\), and
\[
\begin{aligned}
 4^1&\equiv4,& 4^2&\equiv16,& 4^3&\equiv10,\\
 4^4&\equiv13,& 4^5&\equiv25,& 4^6&\equiv19,\\
 4^7&\equiv22,& 4^8&\equiv7,& 4^9&\equiv1\pmod{27}.
\end{aligned}
\]
No preceding power is one; therefore
\(\operatorname{ord}_{27}(4)=9\).
\end{proof}

\subsection{Exact carry identity}

For \(r\in\{0,\ldots,8\}\) and \(n\in\mathbb N\), define
\[
 \kappa_r(n):=\left\lfloor\frac{r+n}{9}\right\rfloor,
 \qquad
 r':=[r+n]_9.
\]

\begin{theorem}[Quaternary carry cocycle]
\label{p12-83-c1-s12:thm:u006c-cociclo-d4}
We have
\begin{equation}
 \boxed{
 \kappa_{[4r]_9}(4n)
 =4\kappa_r(n)
 +\left\lfloor\frac{4\widetilde r'}{9}\right\rfloor
 -\left\lfloor\frac{4\widetilde r}{9}\right\rfloor.}
 \label{p12-83-c1-s12:eq:u006c-cociclo-d4}
\end{equation}
The last difference is the coboundary produced by the choice of
representatives \(r\mapsto\widetilde r\).
\end{theorem}

\begin{proof}
Write
\[
 r+n=9\kappa_r(n)+\widetilde r'.
\]
When multiplied by four,
\[
 4r+4n=36\kappa_r(n)+4\widetilde r'.
\]
Separating quotient and residue modulo nine and subtracting the quotient already contained in
\(4\widetilde r\) yields
\eqref{p12-83-c1-s12:eq:u006c-cociclo-d4}.
\end{proof}

\begin{corollary}[Third iteration]
\label{p12-83-c1-s12:cor:u006c-tercera-iteracion}
For every \((b,r)\in\mathcal Q_{12,9}\),
\begin{equation}
 \boxed{D_4^3(b,r)=([4b+7\widetilde r]_{12},r).}
 \label{p12-83-c1-s12:eq:u006c-d4-cubo}
\end{equation}
\end{corollary}

\begin{proof}
By \cref{p12-83-c1-s12:prop:u006c-d4-total},
\[
 \Theta(D_4^3(b,r))
 =[4^3(9b+\widetilde r)]_{108}
 =[64(9b+\widetilde r)]_{108}.
\]
Since \(64=1+7\cdot9\), the residue modulo nine is again \(r\), whereas
the block coordinate becomes
\([4b+7\widetilde r]_{12}\).
\end{proof}

The return of \(r\) is not return of the complete register. The difference
\(7\widetilde r\) preserves the memory of the internal representative traversed.

\subsection{Obstruction in the observable quotient}

In the observable configuration of
\cref{p12-83-c1-s3:def:tpk-configuracion-observable}, we write
\[
 q=(p^\Sigma,p^\Pi,d^\Sigma,d^\Pi,\theta)
 \in\mathcal Q_{\rm obs}.
\]
The positional projection is
\begin{equation}
 \pi_{\rm pos}(q):=(p^\Sigma,p^\Pi)\in X_9^2.
 \label{p12-83-c1-s12:eq:u006c-proyeccion-posicional}
\end{equation}
The temporal signature is
\[
 \tau(q):=\varepsilon_9(r(\theta))\in\{-1,0,+1\}.
\]
We define the oriented observable shift
\begin{equation}
 a(q):=
 \begin{cases}
  (\delta(d^\Sigma),0),&\tau(q)=+1,\\
  (0,\delta(d^\Pi)),&\tau(q)=-1,\\
  (0,0),&\tau(q)=0.
 \end{cases}
 \label{p12-83-c1-s12:eq:u006c-desplazamiento-observable}
\end{equation}
Then
\[
 \pi_{\rm pos}(F_{\rm obs}q)=\pi_{\rm pos}(q)+a(q),
\]
with the address update and clock defined in U005.

\begin{theorem}[Obstruction to the ordinary observable endofunctor]
\label{p12-83-c1-s12:thm:u006c-obstruccion}
There is no endofunctor of the ordinary observable space that simultaneously satisfies:
\begin{enumerate}
 \item multiplying the geometric position by four;
 \item replace each transition by four genuine microtransitions;
 \item preserving sheet and orientation in each microtransition;
 \item sending a complete duodecaphase fiber to a single
 duodecaphase fiber.
\end{enumerate}
The obstruction persiste in a step of umbral with \(a(q)=0\).
\end{theorem}

\begin{proof}
Let
\[
 \mathcal B_b:=\{(b,r):0\leq r\leq8\}.
\]
Replace one elementary iteration with four force
\[
 h(\theta+1)=h(\theta)+4
\]
and, by induction,
\[
 h(\theta)=4\theta+c\pmod{108}.
\]
Take \(0\leq\widetilde c<108\) and write \(\widetilde c=9c_b+c_0\), with \(0\leq c_0<9\). If the whole
fiber \(\mathcal B_b\) were sent to a single fiber, the destination block index
would have to be independent of \(r\); however, it contains
\begin{equation}
 4b+c_b+
 \left\lfloor\frac{4r+c_0}{9}\right\rfloor
 \pmod{12},
 \label{p12-83-c1-s12:eq:u006c-bloque-obstruccion}
\end{equation}
which is not constant in \(r\). Moreover, the neutral transition operator
satisfies \(P_0^{\rm TPK}=I\), distinct from the cyclotomic
trivializations \(P_b\); the ordinary calendar \((+1,-1,0)\) causes
four elementary iterations to visit both arithmetic modes and the
neutral mode.
\end{proof}

\subsection{Minimal refinement by microinstants}

\begin{definition}[Subdivided observable state]
\label{p12-83-c1-s12:def:u006c-estado-subdividido}
Let
\begin{equation}
 \widetilde{\mathcal Q}_4
 :=\mathcal Q_{\rm obs}\times\{0,1,2,3\},
 \qquad
 \widetilde\pi(q,j)
 :=4\pi_{\rm pos}(q)+j\,a(q).
 \label{p12-83-c1-s12:eq:u006c-estado-refinado}
\end{equation}
We define
\begin{equation}
 \widetilde F(q,j):=
 \begin{cases}
  (q,j+1),&0\leq j<3,\\
  (F_{\rm obs}q,0),&j=3,
 \end{cases}
 \qquad
 \iota_4(q):=(q,0).
 \label{p12-83-c1-s12:eq:u006c-dinamica-refinada}
\end{equation}
\end{definition}

\begin{theorem}[Four-step semiconjugacy]
\label{p12-83-c1-s12:thm:u006c-semiconjugacion}
The inclusion \(\iota_4\) satisfies
\begin{equation}
 \boxed{
 \widetilde F^4\circ\iota_4
 =\iota_4\circ F_{\rm obs}.}
 \label{p12-83-c1-s12:eq:u006c-semiconjugacion}
\end{equation}
The intermediate positions are exactly
\[
 4\pi_{\rm pos}(q)+j\,a(q),
 \qquad j=0,1,2,3.
\]
\end{theorem}

\begin{proof}
Starting from \((q,0)\), the first three applications increment only the
microinstant index. The fourth applies \(F_{\rm obs}\) and returns the index to
zero. This gives
\[
 (q,0)\mapsto(q,1)\mapsto(q,2)\mapsto(q,3)
 \mapsto(F_{\rm obs}q,0).
\]
The formula for the positions is the definition of
\(\widetilde\pi\). The refinement makes visible a determined subdivision;
it does not add a dynamic decision.
\end{proof}

\subsection{Subdivided category and cyclotomic elevation}

Let \(\operatorname{sd}_4\mathsf P_{\rm obs}\) be the category obtained by
replacing each generator
\(\gamma:q\to F_{\rm obs}q\) with the chain
\begin{equation}
 (q,0)\to(q,1)\to(q,2)\to(q,3)\to(F_{\rm obs}q,0).
 \label{p12-83-c1-s12:eq:u006c-cadena-subdivision}
\end{equation}
The functor
\[
 \operatorname{sd}_4:
 \mathsf P_{\rm obs}\longrightarrow
 \operatorname{sd}_4\mathsf P_{\rm obs}
\]
sends each generator to that string. The collapse functor
\[
 \operatorname{col}_4:
 \operatorname{sd}_4\mathsf P_{\rm obs}
 \longrightarrow\mathsf P_{\rm obs}
\]
recompose the four microarrows and satisfies
\begin{equation}
 \operatorname{col}_4\circ\operatorname{sd}_4
 =\operatorname{id}_{\mathsf P_{\rm obs}}.
 \label{p12-83-c1-s12:eq:u006c-colapso}
\end{equation}

In a chain that remains in the fiber \(b\), the first three
microarrows act by \(\rho_b\); the fourth incorporates the interfiber transport
if a dodecaphase crossing occurs. By \(\rho_b^3=I\), the ordered product
of the four actions coincides with the cyclotomic morphism of the
original arrow.

\begin{proposition}[Cyclotomic Invariance under Subdivision]
\label{p12-83-c1-s12:prop:u006c-invariancia-ciclotomica}
There exists a typed lift
\[
 \widetilde{\mathfrak F}_{\rm cyc}:
 \operatorname{sd}_4\mathsf P_{\rm obs}
 \longrightarrow
 \mathbf{Rep}_{\mathbb C}(C_3)\times\mathbf B\mathbb N
\]
such that, for every arrow \(\gamma\) of
\(\mathsf P_{\rm obs}\),
\begin{equation}
 \boxed{
 \widetilde{\mathfrak F}_{\rm cyc}
 \bigl(\operatorname{sd}_4\gamma\bigr)
 =
 \bigl(\mathfrak F_{\rm cyc}\circ Q\bigr)(\gamma).}
 \label{p12-83-c1-s12:eq:u006c-invariancia-ciclotomica}
\end{equation}
\end{proposition}

\begin{proof}
The successive powers of
\(\rho_b\) are assigned to the internal microarrows, and transport between the frames of the
source and target fibers is assigned to the fourth. The composition of the first three factors is the
identity because \(\rho_b^3=I\). The remaining factor is exactly the
morphism that \(\mathfrak F_{\rm cyc}\circ Q\) assigns to the
collapsed arrow. The subdivided category is free on those microarrows, so
the assignment extends uniquely.
\end{proof}

\subsection{\(3\) noninterchangeable quaternary operations}

The architecture contains three distinct operations:
\begin{enumerate}
 \item \(\mathcal S_4\), cardinal substitution of objects and paths;
 \item \(D_4\), multiplication by four in
 \(\Theta_{108}\);
 \item \(\operatorname{sd}_4\), functorial subdivision of an observable
 transition.
\end{enumerate}
Their compatibilities are given, respectively, by
\[
 4^3=1+7\cdot9,\qquad
 \Theta D_4=4\Theta,\qquad
 \widetilde F^4\iota_4=\iota_4F_{\rm obs}.
\]
Sharing the integer four converts none into another.

\begin{criterion}[Refutation criteria]
\label{p12-83-c1-s12:crit:u006c-refutacion}
The construction is refuted if any of the identities
\eqref{p12-83-c1-s12:eq:u006c-d4-total},
\eqref{p12-83-c1-s12:eq:u006c-cociclo-d4},
\eqref{p12-83-c1-s12:eq:u006c-d4-cubo},
\eqref{p12-83-c1-s12:eq:u006c-semiconjugacion}, or
\eqref{p12-83-c1-s12:eq:u006c-invariancia-ciclotomica} fails; likewise if an observable
realization satisfies the four conditions of
\cref{p12-83-c1-s12:thm:u006c-obstruccion} without preserving a register equivalent to the
microinstant.
\end{criterion}


% Unidad U006. Lectores, odómetro y calendario; redacción autónoma.

\section{Modular readers, odometer, and dodecaphasic calendar}
\label{p12-83-c1-s13:chap:tpk-lectores-odometro-calendario}

The enriched state of the TPK is not identified with any of its readings.
This unit defines the readers modulo nine, modulo three, and modulo one thousand, the
phase--sector odometer, the predictive quotient of thirty-six states, and the
cyclic extension of one hundred eight positions. Each construction declares what it
publishes and what memory it preserves.

\subsection{Typed scheme of the modular readers of the TPK: domain, codomain, publication, and conserved memory}

\begin{definition}[Operator specification]
\label{p12-83-c1-s13:def:tpk-ficha-operatoria}
A canonical TPK operator is recorded by means of a sextuple
\[
 \mathcal O=(D_{\mathcal O},C_{\mathcal O},f_{\mathcal O},
 I_{\mathcal O},L_{\mathcal O},M_{\mathcal O}),
\]
where \(D_{\mathcal O}\) is the domain, \(C_{\mathcal O}\) the codomain,
\(f_{\mathcal O}\) the rule, \(I_{\mathcal O}\) the preserved invariants,
\(L_{\mathcal O}\) the information published or forgotten, and
\(M_{\mathcal O}\) the memory that remains in the enriched state.
\end{definition}

Two symbols with equal cardinality or equal period do not represent the same
operator unless there exists an equivalence intertwining their domains,
rules, and invariants. This convention will prevent confusing the predictive quotient
\(C_{36}^{\rm pred}\), the boundary \(R_{36}\), and a
cardinal defect of value thirty-six.

\subsection{Nonadic reader of the enriched state: phase, transition and memory-return}

For \(n\geq1\), define
\[
 \rho_9(n)=1+((n-1)\bmod9),
 \qquad
 q_9(n)=\left\lfloor\frac{n-1}{9}\right\rfloor.
\]
Then
\[
 n=9q_9(n)+\rho_9(n).
\]
The card of \(\rho_9\) is
\[
 \bigl(
 \mathbb Z_{>0},\mathcal D_9,\rho_9,
 n\equiv\rho_9(n)\pmod9,
 \text{positive phase},
 q_9(n)
 \bigr).
\]
The positive digital root numerically coincides with \(\rho_9\), but its
semantics is distinct: one is presented as a section of the quotient; the other,
as an iterated sum of digits.

\subsection{Tritic reader of the oriented state: temporal order and carry transport}

The balanced division is written
\[
 n=r_3(n)+3c_3(n),
 \qquad r_3(n)\in\mathsf{TRIT},\quad c_3(n)\in\mathbb Z.
\]
Its card is
\[
 \bigl(
 \mathbb Z,\mathsf{TRIT}\times\mathbb Z,
 n\mapsto(r_3(n),c_3(n)),
 \operatorname{rec}_3(r,c)=r+3c,
 r_3,
 c_3
 \bigr).
\]
The publication may display only \(r_3\), but the TPK preserves
\(c_3\) if a subsequent composition requires reconstruction.

\subsection{Reader positional by blocks in base \(10^3\)}

The Euclidean division defines
\[
 N=1000q_{1000}(N)+r_{1000}(N),
 \qquad 0\leq r_{1000}(N)<1000.
\]
Its card is
\[
 \bigl(
 \mathbb Z,\mathbb Z\times\{0,\ldots,999\},
 N\mapsto(q_{1000},r_{1000}),
 \operatorname{rec}_{1000}(q,r)=1000q+r,
 r_{1000},
 q_{1000}
 \bigr).
\]

\begin{proposition}[Concatenation and residual reading]
\label{p12-83-c1-s13:prop:tpk-concatenacion-base-mil}
If \(u_1,\ldots,u_m\in\{0,\ldots,999\}\) and
\[
 D_m=1000D_{m-1}+u_m,
 \qquad D_0=0,
\]
then
\[
 D_m=\sum_{j=1}^m u_j1000^{m-j}.
\]
Moreover,
\[
 D_m\equiv\sum_{j=1}^m u_j\pmod9
\]
because \(1000\equiv1\pmod9\).
\end{proposition}

\begin{proof}
The first equality is obtained by induction on \(m\). The second follows
by reducing it modulo nine.
\end{proof}

The congruence does not reconstruct the order of the blocks. The genealogical register preserves the sequence \((u_1,\ldots,u_m)\); the residual reader publishes only its sum modulo nine.

\subsection{Odometer category of the TPK: cylindrical objects and composition by truncation}

\begin{definition}[Nonadic odometer--dodecaphasic]
\label{p12-83-c1-s13:def:tpk-odometro-9-12}
The category \(\mathcal O_{9,12}\) has objects
\[
 (b,r)\in C_{12}\times\{0,\ldots,8\}.
\]
An arrow of length \(n\in\mathbb N\) is
\[
 (b,r)\xrightarrow{\ n\ }
 \left(
  b+\left\lfloor\frac{r+n}{9}\right\rfloor,
  r+n\bmod9
 \right).
\]
The composition sums lengths.
\end{definition}

Let
\[
 \kappa_r(n)=\left\lfloor\frac{r+n}{9}\right\rfloor.
\]

\begin{lemma}[Cocycle of the odometer]
\label{p12-83-c1-s13:lem:tpk-cociclo-odometro}
If \(r_1=r+n_1\bmod9\), then
\[
 \kappa_r(n_1+n_2)
 =\kappa_r(n_1)+\kappa_{r_1}(n_2).
\]
\end{lemma}

\begin{proof}
We write \(r+n_1=9\kappa_r(n_1)+r_1\), with
\(0\leq r_1<9\). Adding \(n_2\) and again dividing by nine
yields the identity.
\end{proof}

For a state \(x\) whose base contains \(\theta\), we define
\[
 Q_0(x)=(\beta(\theta),\widetilde\theta\bmod9).
\]
A TPK arrow has length equal to the number of observable generators that
compose it.

\begin{theorem}[Phase--length functor]
\label{p12-83-c1-s13:thm:tpk-funtor-fase-longitud}
The assignment sending a state \(x\) to \(Q_0(x)\) and an arrow of
length \(n\) to the odometric arrow of length \(n\) defines a functor
\[
 Q:\mathsf{TPK}\longrightarrow\mathcal O_{9,12}.
\]
\end{theorem}

\begin{proof}
Each generator increases \(\theta\) by one unit. After \(n\) steps, the
phase residue increases by \(n\) modulo nine and the sector increases by
\(\lfloor(r+n)/9\rfloor\). The identity has zero length and the
\cref{p12-83-c1-s13:lem:tpk-cociclo-odometro} proves compatibility with composition.
\end{proof}

After nine steps,
\[
 (b,r)\longmapsto(b+1,r).
\]
After one hundred and eight,
\[
 (b,r)\longmapsto(b,r).
\]
These identities belong to the odometric quotient. They do not assert that the
cylinder, boundary, genealogical register, or vertical memory have
returned.

\subsection{Minimum predictive-memory quotient}

Let
\[
 X_{9,12}=\{(a,m):a\in\mathbb Z/9\mathbb Z, 0\leq m<12\}
\]
with successor
\[
 S(a,m)=
 \begin{cases}
  (a,m+1),&m<11,\\
  (a+1\bmod9,0),&m=11.
 \end{cases}
\]
We define the observables
\[
 \beta_{\rm pred}(a,m)=4a\pmod{12},
 \qquad
 d_{\rm pred}(a,m)=1+(m+\beta_{\rm pred}(a,m)\pmod{12}).
\]
For an observable \(q\), the inclusive future word of horizon \(h\)
is
\[
 Q_h^q(x)=(q(x),q(Sx),\ldots,q(S^hx)).
\]

\begin{theorem}[Minimum predictive quotient]
\label{p12-83-c1-s13:thm:tpk-cociente-predictivo-36}
The observables \(\beta_{\rm pred}\), \(d_{\rm pred}\), and the pair
\((\beta_{\rm pred},d_{\rm pred})\) produce the same stable quotient
\[
 \pi_{36}(a,m)=(a\bmod3,m).
\]
It has thirty-six states and fibers of cardinality three. The minimal
horizons are \(11\), \(8\), and \(0\), respectively.
\end{theorem}

\begin{proof}
For a finite system \((X,T)\), let
\[
 E_h=\operatorname{Eq}(Q_h^q).
\]
Then
\[
 E_{h+1}=E_h\cap(T\times T)^{-1}(E_h).
\]
The decreasing chain stabilizes and its limit is the largest
\(T\)-compatible equivalence contained in \(\operatorname{Eq}(q)\). In the declared
system, the exact enumeration produces the following number of classes:
\[
\begin{array}{c|rrrrrrrrrrrr}
h&0&1&2&3&4&5&6&7&8&9&10&11\\\hline
\#(X/E_h^{\beta})&3&6&9&12&15&18&21&24&27&30&33&36
\end{array},
\]
and
\[
\begin{array}{c|rrrrrrrrr}
h&0&1&2&3&4&5&6&7&8\\\hline
\#(X/E_h^{d})&12&15&18&21&24&27&30&33&36.
\end{array}
\]
The pair already distinguishes the thirty-six classes at horizon zero. In all
cases, the stable classes preserve \(m\) and \(a\bmod3\); the three
possibilities for \(a\) with the same residue modulo three form each fiber.
\end{proof}

We denote this object by \(C_{36}^{\rm pred}\). No is the boundary
\(R_{36}\), that will contain three ordered channels of six trits, nor a census
of thirty-six elements of another construction.

\subsection{Central extension of the calendar}

We define
\[
 \iota:C_{12}\longrightarrow C_{108},
 \qquad \iota([b]_{12})=[9b]_{108},
\]
and
\[
 p:C_{108}\longrightarrow C_9,
 \qquad p([\theta]_{108})=[\theta]_9.
\]
The exact sequence is obtained.
\begin{equation}
 0\longrightarrow C_{12}
 \xrightarrow{\ \iota\ }C_{108}
 \xrightarrow{\ p\ }C_9
 \longrightarrow0.
\label{p12-83-c1-s13:eq:tpk-extension-12-108-9}
\end{equation}
A section of sets \(s([r]_9)=[r]_{108}\),
\(0\leq r<9\), induces the cocycle
\[
 c(r,r')=
 \left[\left\lfloor\frac{r+r'}{9}\right\rfloor\right]_{12}.
\]

\begin{proposition}[The extension does not split]
\label{p12-83-c1-s13:prop:tpk-extension-no-escindida}
The sucesion \cref{p12-83-c1-s13:eq:tpk-extension-12-108-9} not is isomorfa to the extension
product \(C_{12}\times C_9\).
\end{proposition}

\begin{proof}
The exponent of \(C_{108}\) is \(108\). The exponent of
\(C_{12}\times C_9\) is
\(\operatorname{mcm}(12,9)=36\). Two isomorphic finite groups have the
same exponent; therefore, they are not isomorphic.
\end{proof}

\begin{theorem}[Necessity and exact period of one hundred eight positions]
\label{p12-83-c1-s13:thm:tpk-periodo-exacto-108}
The translation \(\theta\mapsto\theta+1\) on \(C_{108}\) has order
exactly \(108\). Consequently, an iteration \(n\) returns the
complete calendar if and only if \(108\mid n\).
\end{theorem}

\begin{proof}
The \(n\)-th power sends \([\theta]\) to \([\theta+n]_{108}\). It is
the identity for every \(\theta\) exactly when \([n]_{108}=0\), that is,
when \(108\mid n\).
\end{proof}

In particular, \(1008\) is not another return of the calendar:
\[
 1008\equiv36\pmod{108}.
\]
After \(1008\) steps, the phase modulo nine returns, but the coordinate of
\(C_{108}\) has shifted by thirty-six positions. Naturally,
\(216,324,\ldots\) also return the calendar, but they are repetitions of
the fundamental circuit of length \(108\).

\subsection{Unbounded vertical memory}

Let \(q_{\rm mem}\in\mathbb Z_{\geq0}\) be the unbounded number of nonadic
turns, and let
\[
 \beta=[q_{\rm mem}]_{12}.
\]
Then a nine-step turn satisfies
\[
 q_{\rm mem}\longmapsto q_{\rm mem}+1,
 \qquad
 \beta\longmapsto\beta+1\pmod{12}.
\]
After one hundred eight steps,
\[
 \beta\longmapsto\beta,
 \qquad
 q_{\rm mem}\longmapsto q_{\rm mem}+12.
\]
The calendar returns; the vertical memory does not.

\begin{corollary}[Hierarchy of returns]
\label{p12-83-c1-s13:cor:tpk-jerarquia-retornos}
In the phase quotient,
\(9\) steps constitute one turn. The composite morphisms
\[
 \mathcal K_{27}=F_{\rm obs}^{27},
 \qquad
 F_{\rm obs}^{54},
 \qquad
 F_{\rm obs}^{108}
\]
contain, respectively, three, six, and twelve nonadic turns. None
erases the genealogical record or the boundary by definition.
\end{corollary}

\subsection{Minimal type system of the enriched state and compatibilities between its readers}

The following table answers explicitly which internal objects
already exist in the TPK before introducing seeds:
\[
\begin{array}{p{.23\textwidth}p{.27\textwidth}p{.38\textwidth}}
\toprule
Objeto u operador&Dominio y codominio&Función\\\midrule
\(X_9\)&\((\mathbb Z/9)^2\)&soporte APP\\
\(T_N,T_E,T_S,T_O\)&\(X_9\to X_9\)&traslaciones cardinales\\
\(M_{\rm ph}\)&\(\mathcal D_9\to\mathsf{TRIT}\)&selección de hoja\\
\(F_{\rm obs}\)&\(\mathcal Q_{\rm obs}\to\mathcal Q_{\rm obs}\)&sucesor observable\\
\((\rho_9,q_9)\)&\(\mathbb Z_{>0}\to\mathcal D_9\times\mathbb Z_{\ge0}\)&fase y altura\\
\((r_3,c_3)\)&\(\mathbb Z\to\mathsf{TRIT}\times\mathbb Z\)&orientación y carry\\
\((q_{1000},r_{1000})\)&\(\mathbb Z\to\mathbb Z\times\{0,\ldots,999\}\)&bloque decimal y cociente\\
\(Q\)&\(\mathsf{TPK}\to\mathcal O_{9,12}\)&fase y longitud\\
\(C_{36}^{\rm pred}\)&cociente de \(X_{9,12}\)&memoria predictiva mínima\\
\(C_{108}\)&calendario cíclico&extensión no escindida\\
\(q_{\rm mem}\)&\(\mathbb Z_{\ge0}\)&memoria vertical no acotada\\
\(\operatorname{Ext}_r\)&
\(\operatorname{Ext}_r\subseteq\mathcal F_q\times\mathcal F_{q'}\)&
prolongaciones compatibles\\
\bottomrule
\end{array}
\]

\subsection{Obstructions of the readers and the TPK calendar: Euclidean quotients, \(C_{36}^{\mathrm{pred}}\), extension \(C_{108}\), and vertical memory}

\begin{criterion}[Reader Refutation Criteria and Calendar]
The unit is refuted if:
\begin{enumerate}
 \item a residual reader is presented as a complete transition;
 \item the quotient of some subsequently used Euclidean division is omitted;
 \item \(C_{36}^{\rm pred}\) is identified with \(R_{36}\);
 \item it is asserted that the extension \(C_{108}\) is the product
       \(C_{12}\times C_9\);
 \item complete return is declared for a number not divisible by \(108\);
 \item the return of \(\beta\) fuerza \(q_{\rm mem}=0\);
 \item \(\mathcal K_{27}\) is presented as a distinct clock and not as
       a power of the transport.
\end{enumerate}
\end{criterion}

\section{Complete operatorial structure of APP and arithmetic-geometric hierarchy}

The operatorial development of APP fixes the hierarchy
\(4\to2\to6\to54\), the arithmetic projectors, the central block, and the
local geometry of paths. These consequences remain subordinated to the
already constructed formal kernel and do not introduce a second origin of APP.

The next unit will construct the quadruple subdivision, the
cyclotomic functor, and the internal channels of ninety and one hundred twenty positions.
Only afterward will the exhaustive space of seeds be introduced.
